\documentclass[a4paper,12pt]{book}

\usepackage[unicode,colorlinks=true,linkcolor=blue]{hyperref}
% здесь подключении шрифтов в русскими буквами
\usepackage[T2A]{fontenc}
\usepackage[utf8]{inputenc}
\usepackage[english,russian]{babel}
\usepackage{amsmath,amsthm,amssymb,amsfonts,mathtext,cite,enumerate,float}
%\usepackage[dvips]{graphicx}
\usepackage[pdftex]{graphicx}
\graphicspath{{images/}}
\usepackage{fix-cm}
\usepackage{mathenv}
\usepackage{arcs}
\usepackage{yhmath}
\usepackage{indentfirst}

\makeatletter
\newcommand*{\rom}[1]{\expandafter\@slowromancap\romannumeral #1@}
\renewcommand{\@biblabel}[1]{#1.}
\makeatother

\usepackage{geometry}  % Меняем поля страницы
\geometry{left=2cm}    % левое поле
\geometry{right=1.5cm} % правое поле
\geometry{top=2cm}     % верхнее поле
\geometry{bottom=2cm}  % нижнее поле

\newcommand{\HRule}{\rule{\linewidth}{0.5mm}}

\renewcommand{\theenumi}{\arabic{enumi}}                                       % Меняем везде перечисления на цифра.цифра
\renewcommand{\labelenumi}{\arabic{enumi}.}                                     % Меняем везде перечисления на цифра.цифра
\renewcommand{\theenumii}{.\arabic{enumii}}                                    % Меняем везде перечисления на цифра.цифра
\renewcommand{\labelenumii}{\arabic{enumi}.\arabic{enumii}.}                   % Меняем везде перечисления на цифра.цифра
\renewcommand{\theenumiii}{.\arabic{enumiii}}                                  % Меняем везде перечисления на цифра.цифра
\renewcommand{\labelenumiii}{\arabic{enumi}.\arabic{enumii}.\arabic{enumiii}.} % Меняем везде перечисления на цифра.цифра

\theoremstyle{plain}
\newtheorem{theorem}{Теорема}[section]
\newtheorem*{consequence}{Следствие}
\newtheorem{lemma}{Лемма}[section]
\theoremstyle{definition}
\newtheorem{definition}{Определение}[section]
\theoremstyle{remark}
\newtheorem*{comment}{Комментарий}
\newtheorem*{clarification}{Пояснение}
\newtheorem{example}{Пример}[section]
\newtheorem{exercise}{Упражнение}[]
\newtheorem{remark}{Замечание}[section]
\newtheorem{approval}{Утверждение}[section]

\begin{document}
  \begin{titlepage}
  \begin{center}

    % Upper part of the page. The '~' is needed because \\
    % only works if a paragraph has started.
    %\includegraphics[width=0.15\textwidth]{./logo}~\\[1cm]

    \textsc{\LARGE НИЯУ МИФИ}\\[1.5cm]

    \textsc{\Large 3 семестр факультет КиБ}\\[0.5cm]

    % Title
    \HRule \\[0.4cm]
    {\huge \bfseries Математический анализ\\[0.4cm]}

    \HRule \\[1.5cm]

    \begin{minipage}{0.4\textwidth}
      \begin{flushleft} \large
        \emph{Автор:}\\
        Тропин \textsc{А.Г.}
      \end{flushleft}
    \end{minipage}
    \begin{minipage}{0.4\textwidth}
      \begin{flushright} \large
        \emph{Лектор:} \\
        Теляковский \textsc{Д.С.}
      \end{flushright}
    \end{minipage}

    \vfill
    \begin{flushleft}
      e-mail: \href{mailto:andrewtropin@gmail.com}{andrewtropin@gmail.com} \\
      github: \href{http://github.com/abcdw/mephi}{abcdw/mephi}
    \end{flushleft}
    {\large \today}
  \end{center}
\end{titlepage}

  \tableofcontents

  \part{Функциональные последовательности и ряды}
  \chapter{Числовые ряды}
  \section{Определение}

\begin{definition}
  $U_1 + U_2 + U_3 + \dots = \sum\limits_{k = 0}^{\infty}{U_k}$
\end{definition}

\begin{definition}[Частичная сумма]
  $S_n = \sum\limits_{k = 0}^n U_k$
\end{definition}

\begin{definition}
  Ряд сходится, если $\exists \lim\limits_{n \to \infty}
                      \sum\limits_0^\infty U_k = S$
\end{definition}

\begin{definition}
  $\{a_n\} = a_0 + \sum\limits_1^n (a_k - a_{k-1})$, где $(a_k - a_{k - 1}) = U_k$
\end{definition}

\begin{theorem}[Критерий Коши]
  \label{th111}
  Ряд сходится, тогда и только тогда, когда он удовлетворяет условию Коши: \\
  $$\forall \varepsilon > 0, \ \exists N~=~N(\varepsilon), \
  \forall n \geq N, \forall p: \ |\sum\limits_{k=n+1}^{n+p} U_k|~=~|U_{n+1}
  + \dots + U_{n+p}|~=~|S_{n+p} - S_{n}| < \varepsilon$$
\end{theorem}

\begin{proof}
  $\sum U_k$ - сходится $\Leftrightarrow \{S_n\}$
  $$\forall \varepsilon > 0, \ \exists N: \forall n \geq N, \forall p: \
  |S_{n+p} - S_{n}| < \varepsilon$$
\end{proof}
\begin{consequence}[Необходимое условие сходимости]
  Если $\sum U_k$ сходится, то $U_k \to 0$, при $k~\to~\infty$
\end{consequence}
\begin{proof}
  Если $\sum U_k$ сходится, то выполняется Критерий Коши. При $p = 1$
  $$\forall \varepsilon > 0, \exists N, \forall n \geq N:
  |S_{n+1} - S_n|~=~|U_{n+1}| < \varepsilon$$
\end{proof}
\begin{consequence}
  Отбрасывание или добавление любого конечного числа членов ряда на его
  сходимость не влияет.
\end{consequence}

\begin{example}
  $\sum\limits_0^\infty z^n,$
  $S_n(z) = \sum\limits_0^n z_n = \frac{1-z^{n+1}}{1-z}$.
  При $n \to \infty, \ S_n(z) = \frac{1}{1 - z}, \ |z| < 1$. \\
  $S_n(z)$ не имеет придела при $|z| \geq 1$.
\end{example}

\section{Действия с рядами}

\begin{theorem}
  Ряды $\sum U_k$ и $\sum V_k$ сходятся, $\alpha$ --- комплексное число, тогда
  \begin{align}
    \label{series:property1}
    \sum\alpha U_k &= \alpha\sum U_k \\
    \label{series:property2}
    \sum (U_k \pm V_k) &= \sum U_k \pm \sum V_k
  \end{align}
\end{theorem}
\begin{proof}
  Доказательство свойства (\eqref{series:property1}):
  $$\sum\limits_{k=0}^\infty \alpha U_k = \lim\limits_{n \to \infty}
  \sum\limits_0^n \alpha U_k = \alpha \lim\limits_{n \to \infty} \sum\limits_0^n U_k
  = \alpha \sum\limits_{k = 0}^\infty U_k$$
\end{proof}
\begin{proof}
  Доказательство свойства (\eqref{series:property2}):
  \begin{gather*}
      \sum\limits_0^\infty U_k \pm \sum\limits_{0}^{\infty} V_k
      = \lim\limits_{n \to \infty} \sum\limits_{0}^{n} U_k
      \pm \lim\limits_{n \to \infty} \sum\limits_{0}^{n} V_k = \\
      = \lim\limits_{n \to \infty} \sum\limits_{0}^{n} U_k
      \pm \sum\limits_{0}^{n} V_k
      = \lim\limits_{n \to \infty} \sum\limits_{0}^{n} (U_k \pm V_k)
      = \sum\limits_{0}^{n} (U_k \pm V_k)
  \end{gather*}
\end{proof}

\begin{remark}
  Из сходимости $\sum (U_k \pm V_k) \not \Rightarrow$ сходимость
  $\sum U_k$ и $\sum V_k$
\end{remark}
\begin{remark}
  Если $\sum U_k$ сходится, то можно группировать, не меняя порядка.
\end{remark}
\begin{example}
  \begin{gather*}
    \sum (1 - 1) \\
    (1 - 1) + (1 - 1) + \dots \\
    1 - (1 - 1) - (1 - 1) \dots
  \end{gather*}
\end{example}
\begin{comment}
  Нельзя раскрывать скобки и переставлять члены.
\end{comment}
\subsection{Ряды с неотрицательными членами}
$U_k \geq 0$, $S_n = \sum\limits_{0}^{n} U_k$ - не убывающая последовательность. \\
$\sum\limits_{0}^{n} U_k$ - сходится $\Leftrightarrow$ $\{S_n\} $ --- ограничена
\begin{comment}
  Сходимость ряда эквивалентна ограниченности $S_n$
\end{comment}

\begin{theorem}
    $$U_k \geq 0, V_k \geq 0, \ \forall k:$$
    \begin{enumerate}
      \item
        Если $0 \leq U_k \leq V_k$, то если $\sum V_k$ сходится $\Rightarrow$
        $\sum U_k$ сходится и если
        $\sum U_k$ расходится $\Rightarrow$ $\sum V_k$ расходится.
      \item
        Если $\lim\limits_{n\to \infty} \frac{U_k}{V_k} = A > 0$, то ряды
        сходятся или расходятся одновременно.
    \end{enumerate}
\end{theorem}

\pagebreak
\begin{proof}
  \hfill
  \begin{enumerate}
    \item
      $\forall n$ верно неравенство $0 \leq \sum\limits_{0}^{n} U_k \leq
      \sum\limits_{0}^{n} V_k$
    \item
      $\forall \varepsilon > 0 \ \varepsilon < A \ \exists N : \forall n \geq N
      \Rightarrow 0 < A - \varepsilon < \frac{U_k}{V_k} < A + \varepsilon$ \\
      $0 < (A - \varepsilon) \cdot V_k < U_k < (A + \varepsilon) \cdot V_k$ \\
      Пусть $U_k$ --- сходится, тогда из доказанного выше 1ого пункта следует
      $(A - \varepsilon) \cdot V_k$ --- сходится
      $\Rightarrow \sum V_k$ сходится
      $\Rightarrow \sum (A + \varepsilon) \cdot V_k$ сходится
      $\Rightarrow \sum U_k$ сходится.
  \end{enumerate}
\end{proof}

\begin{remark}
  Вместо существования предела $\lim\limits_{n \to \infty} \frac{U_k}{V_k}$
  достаточно предположить, что существуют такие числа
  p и q > 0, такие что $0 < q < \frac{U_k}{V_k} < p, \ \forall k$
\end{remark}

\begin{theorem}[Признак Даламбера]
  $$\sum U_k, \ U_k > 0$$
  \begin{enumerate}
    \item Если $\exists q$ такое что:
      $\ \forall k \ \frac{U_{k+1}}{U_k} \leq q < 1$ сходится
    \item Если $\exists \lim\limits_{k \to \infty} \frac{U_{k+1}}{U_k} = q$, то:
    \begin{itemize}
      \item при $q < 1$ сходимость
      \item при $q > 1$ расходимость
      \item при $q = 1$ неизвестно (нужно провести дополнительные исследования)
    \end{itemize}
  \end{enumerate}
\end{theorem}

\begin{proof}
  Идея докозательства - сравнение с геометрической прогрессией.
  \begin{enumerate}
    \item $k = 0, 1, \dots, n$;
      $U_k = U_0 \cdot \frac{U_1}{U_0} \frac{U_2}{U_1} \cdots \frac{U_k}{U_{k-1}}
      < U_0 \cdot q^k$
      \begin{comment}
        $\frac{U_k}{U_{k-1}} < q, \ \forall k$
      \end{comment}
      $q < 1$, тогда $\sum U_0 \cdot q^k$ --- сходящаяся геометрическая прогрессия. \\
      $U_k = U_0 \cdot \frac{U_1}{U_0}\frac{U_2}{U_1}\cdots\frac{U_k}{U_{k-1}}
      \geq U_0 > 0$
      \begin{comment}
        $\frac{U_k}{U_{k-1}} \geq 1, \ \forall k$
      \end{comment}
      $U_k \not \to 0 \Rightarrow$ не выполняется необходимое условие сходимости.
    \item Пусть $\lim\limits_{k \to \infty} \frac{U_{k+1}}{U_k} = q$ \\
      $\forall \varepsilon > 0, \ \exists K : \ \forall k \geq K$
      выполняется неравенство
      $q - \varepsilon < \frac{U_{k+1}}{U_k} < q + \varepsilon$
      \begin{itemize}
        \item Если $q < 1$,выберем такое $\varepsilon$, что 
          $q + \varepsilon < 1$, для
          $\forall k \geq K(\varepsilon)$. \\
          $\frac{U_{k+1}}{U_k} < q + \varepsilon < 1 \Rightarrow$
          сходится по первой части.
        \item Если $q > 1$, то выберем $\varepsilon$ так, чтобы
          $q - \varepsilon > 1$, для
          $\forall k \geq K(\varepsilon)$. \\
          $\frac{U_{k+1}}{U_k} > q - \varepsilon > 1, \ \Rightarrow$
          расходится по первой части.
      \end{itemize}
    \end{enumerate}
\end{proof}

\pagebreak
\begin{theorem}[Признак Коши]
  $$\sum U_k, U_k \geq 0$$
  \begin{enumerate}
    \item Если $\exists q < 1$ и $\forall k > K : $ выполняется
      $\sqrt[k]U_k \leq q < 1$, то ряд сходится, а если
      $\forall k \ \sqrt[k]U_k~\geq~1$, то расходится.
    \item Если $\exists \lim\limits_{k \to \infty} \sqrt[k]U_k = q, (q \geq 0)$,
      то
      \begin{itemize}
        \item $q < 1$ - сходится
        \item $q > 1$ - расходится
        \item $q = 1$ - нужны дополнительные исследования
      \end{itemize}
  \end{enumerate}
\end{theorem}
\begin{remark}
  $\lim\limits_{n \to \infty} \sqrt[n]U_n$ можно рассматривать вместо
  $\overline {\lim\limits_{k \to \infty}} \sqrt[k]U_k$
\end{remark}

\begin{proof}
  Сравнение с геометрической прогрессией
  \begin{enumerate}
    \item Если $\forall k\ \sqrt[k]U_k \leq q < 1 \Rightarrow U_k \leq q^k$
      --- сходящаяся геометрическая прогрессия. \\
      Если $\forall k \sqrt[k]U_k \geq 1 \Rightarrow U_k \geq 1$
      --- не выполняется необходимое условие сходимости.
    \item Если $\lim\limits_{k \to \infty} \sqrt[k]U_k = q$, то
      $\forall \varepsilon > 0 \ \exists K = K(\varepsilon) :
      \forall k \geq K, (q - \varepsilon) < \sqrt[k]U_k < (q + \varepsilon)$\\
      $(q - \varepsilon)^k < U_k < (q + \varepsilon)^k$
      \begin{itemize}
        \item При $q < 1$ выберем $\varepsilon$ так, чтобы $q + \varepsilon < 1$,
          тогда $U_k < (q + \varepsilon)^k < 1$
          --- сходящаяся геометрическая прогрессия.
        \item При $q > 1$ выберем $\varepsilon$ так, чтобы $q - \varepsilon > 1$,
          тогда $U_k > (q - \varepsilon)^k > 1$
          --- не выполняется необходимое условие сходимости.
      \end{itemize}
  \end{enumerate}
\end{proof}

\begin{definition}
  Дана $\{a_n\}$ и пусть $\overline {\lim\limits_{n\to \infty}} a_n$
  --- наибольший из частичных пределов, тогда:
  $$\forall \{a\} \ \exists \overline{\lim\limits_{n \to \infty}} a_n = A \ or \ \infty $$
  \begin{comment}
    $A$ --- число.
  \end{comment}
  \begin{itemize}
    \item Если $\overline{\lim\limits_{n \to \infty}} a_n = +\infty \Rightarrow
      \{a_n\}$ --- неограничена сверху $\Rightarrow
      \overline{\lim\limits_{k \to \infty}} \sqrt[k]U_k = +\infty$
      --- неограничена сверху. $U_k$ неограничена сверху и не выполняется
      необходимое условие.
    \item Если $\overline{\lim\limits_{n \to \infty}} a_n = A$, тогда
      $\forall \varepsilon \in (A - \varepsilon, A + \varepsilon)$ бесконечно
      много членов $\{a_n\}$: \\
      \begin{itemize}
        \item $\overline{\lim\limits_{k \to \infty}} \sqrt[k]U_k = q < 1$.
          Выберем $\varepsilon$ так, чтобы $q + \varepsilon < 1 \Rightarrow$
          $\exists K: \forall k \geq K, \ \sqrt[k]U_k < q + \varepsilon < 1$ по признаку
          Коши. \\
        \item $\overline{\lim\limits_{k \to \infty}} \sqrt[k]U_k = q > 1$.
          Выберем $\varepsilon$ так, чтобы $q - \varepsilon > 1 \Rightarrow$
          $\forall K \ \exists k \geq K: \sqrt[k]U_k~>~q~-~\varepsilon~>~1 \\
          \Rightarrow~U_k~>~1$
      \end{itemize}
  \end{itemize}
\end{definition}

\section{Интегральный признак сходимости рядов с неотрицательными членами}
\begin{theorem}
  Если $f(x)$ не отрицательна и убывает на $x \geq 1$, то ряд
  \begin{gather}
    \label{th131:series1}
    \sum\limits_{n = 1}^{\infty} f(n) < \infty
  \end{gather}
  сходится тогда и только тогда, когда сходится интеграл:
  \begin{gather}
    \label{th131:integral1}
    \int\limits_1^{+\infty} f(x) dx
  \end{gather}
  то есть $ \int\limits_1^{+\infty} f(x) dx < \infty$. \\
  \begin{itemize}
    \item $\sum a_n < \infty$ --- сходится
    \item $\sum a_n = \infty$ --- расходится
  \end{itemize}
\end{theorem}

\begin{proof}
  Если $k \leq x \leq k + 1, \ k = 1, 2, \dots,$ то, в силу убывания функции
  получаем неравенство:
  $$f(k) \geq f(x) \geq f(k + 1)$$
  Интегрируя по отрезку $[k, k + 1]$ получим:
  $$f(k) \geq \int\limits_k^{k+1} f(x) dx \geq f(k+1), \ k = 1, 2, \dots \ .$$

  \begin{gather}
    \label{th131:series2}
    \sum\limits_{k = 1}^{n} f(k+1) \leq \int\limits_1^{n+1} f(x) dx \leq
    \sum\limits_{k = 1}^{n} f(k)
  \end{gather}
  Пусть $S_n = \sum\limits_{k = 1}^{n} f(k),$ тогда \eqref{th131:series2} примет вид:
  \begin{gather}
    \label{th131:series3}
    S_n + 1 - f(1) \leq \int\limits_1^{n+1} f(x) dx \leq S_n
  \end{gather}
  Если ряд \eqref{th131:series1} сходится и его сумма равна $S$, то $S_n \leq S$,
  и $\int\limits_1^{n+1} f(x) dx \leq S, \ \forall n \in \mathbb{N}$. \\
  $\forall b > 1, n + 1 > b$ имеем:
  $$\int\limits_1^b f(x) dx \leq \int\limits_1^{n+1} f(x) dx \leq S$$
  В силу неотрицательности функции $f(x)$ интеграл сходится. \\
  Пусть наоборот, интеграл \eqref{th131:integral1} сходится, тогда из
  \eqref{th131:series3} следует:
  $$S_{n+1} \leq f(1) + \int\limits_1^{n+1} f(x) dx
  \leq f(1) + \int\limits_1^\infty f(x) dx$$
  Тем самым, последовательность сумм $\{S_n\}$ ряда \eqref{th131:series1}
  ограничена сверху, и поэтому этот ряд сходится.
\end{proof}

\begin{example}
  \label{th131:example131}
  \begin{gather}
    \label{ex131:series1}
    \sum\limits_{n = 1}^{\infty} \frac{1}{n^\alpha}, \alpha \in \mathbb{R}
  \end{gather}
  Положим $f(x) = \frac{1}{x^\alpha}$, тогда $f(n) = \frac{1}{n^\alpha}$ \\
  Поскольку $\int\limits_1^{+\infty} \frac{dx}{x^\alpha}$:
  \begin{itemize}
    \item При $\alpha > 1$ сходится
    \item При $\alpha \leq 1$ расходится
  \end{itemize}
  Тогда ряд \eqref{ex131:series1} сходится тогда и только тогда, когда $\alpha > 1$.
  При $\alpha < 0 \ $ дробь $ \ \frac{1}{n^\alpha} \geq 1$.
\end{example}

\section{Признак сходимости для чередующихся рядов}
Рассмотрим ряды с действительными числами, которые то положительные, то
отрицательные.

\begin{theorem}[Лейбница]
  Если
  \begin{gather}
    \lim\limits_{n \to \infty} U_n = 0 \label{th141:lim1} \\
    U_n \geq U_{n+1} > 0, \ n = 1, 2, \dots \label{th141:ineq1}
  \end{gather}
  то знакочередеющийся ряд
  \begin{gather}
    \sum\limits_{n = 1}^{+\infty} {(-1)}^{n+1} U_n \label{th141:sum1}
  \end{gather}
  сходится, при этом если $S$ --- сумма ряда, а $S_n$ --- его $n$-ая частичная
  сумма,\\ то $\forall n: n = 1, 2, \dots$
  \begin{gather}
    |S - S_n| \leq U_{n+1} \label{th141:ineq2}
  \end{gather}
\end{theorem}

\begin{proof}
  Заметим, что частичная суммы $S_n$ с четными номерами возрастают:
  $$S_{2k} = (U_1 - U_2) + (U_3 - U_4) + \dots + (U_{2k - 1} - U_{2k}), \
  k = 1, 2, \dots $$
  Так что выполняется неравенство $S_{2k + 2} \geq S_k$. Кроме того, они
  ограничены сверху:
  $$S_{2k} = U_1 - (U_2 - U_3) - \dots - (U_{2k - 2} - U_{2k - 1}) - U_{2k}, \
  S_{2k} < U_1$$
  Поэтому последовательность $\{S_{2k}\}$ сходится
  \begin{gather}
    \lim\limits_{k \to \infty} S_{2k} = S \label{th141:lim2}
  \end{gather}
  Поскольку $S_{2k + 1} = S_{2k} + U_{2k + 1}$ и $U_{2k + 1} \to 0$ при
  $k \to \infty$, то
  \begin{gather}
    \lim\limits_{k \to \infty} S_{2k + 1} = S \label{th141:lim3}
  \end{gather}
  Из \eqref{th141:lim2} и \eqref{th141:lim3} следует, что
  $\lim\limits_{n \to \infty} S_n = S$ \\
  При этом, нетрудно увидеть, что
  \begin{gather}
    S_{2k} \leq S \leq S_{2k + 1} \leq S_{2k - 1}, \ \forall k \label{th141:ineq3}
  \end{gather}
  Из неравенства \eqref{th141:ineq3} следует, что
  \begin{align*}
    S - S_{2k} \leq S_{2k + 1} - S_{2k} &= U_{2k + 1} \\
    S_{2k - 1} - S \leq S_{2k - 1} - S_{2k} &= U_{2k}, \ k = 1, 2, \dots
  \end{align*}
  Это и означает, что $\forall n \in \mathbb{N}$ выполняется
  неравенство \eqref{th141:ineq2}.
\end{proof}

\section{Преобразование Абеля}
\begin{theorem}
  Пусть $a_k \in \mathbb{C}$, $b_k \in \mathbb{C}$, $k = 1, \dots, n$;
  $B_k = b_1 + \dots + b_k$, тогда
  \begin{gather}
    \sum\limits_{k = 1}^{n} a_k b_k
  = \sum\limits_{k = 1}^{n - 1} (a_k - a_{k+1}) b_k + a_n B_n
  \label{th151:sum1}
  \end{gather}
\end{theorem}

\begin{proof}
  Очевидно, $b_1 = B_1, b_k = B_k - B_{k - 1}, k = 2, 3, \dots, n$ \\
  Поэтому
  $a_1 b_1 + a_2 b_2 + \dots + a_n b_n
  = a_1 B_1 + a_2 (B_2 - B_1) + a_3 (B_3 - B_2) + \dots a_n (B_n - B_{n - 1})
  = (a_1 - a_2) B_1 + (a_2 - a_3) B_2 + \dots + (a_{n-1} - a_n) B_{n-1} + a_n B_n$ \\
  Называется преобразованием Абеля $\sum\limits_{k = 1}^{n} a_k b_k$.
\end{proof}

\begin{consequence}[лемма Абеля]
  \label{th151:cons}
  Если $a_1 \leq a_2 \leq \dots \leq a_n$
   или $a_1 \geq a_2 \geq \dots \geq a_n$ \\
  $a_k \in \mathbb{R}, \ \forall k = 1, 2, \dots, n$,
  $|b_1 + \dots + b_k| \leq B, \ (b_k \in \mathbb{C})$, то
  $$\left|\sum\limits_{k = 1}^{n} a_k b_k \right| \leq B(|a_1| + 2|a_n|)$$
\end{consequence}
\begin{proof}
  $\left|\sum\limits_{k = 1}^{n} a_k b_k \right | \leq
  \sum\limits_{k = 1}^{n - 1} |a_k - a_{k + 1}| |B_k| + |a_n B_n| \leq
  B\left(\sum\limits_{k = 1}^{n - 1} |a_k - a_{k + 1}| + |a_n|\right) = \\
= B\left(\left| \sum\limits_{k = 1}^{n - 1} (a_k - a_{k + 1}) \right | + |a_n| \right)
= B(|a_1 - a_n| + |a_n|) \leq B(|a_1| + 2|a_n|).$
\end{proof}

\section{Признаки Дирихле и Абеля}
\begin{theorem}[признак Дирихле]
  \label{th161}
  Пусть дан ряд
  \begin{gather}
    \sum\limits_{n = 1}^{\infty} a_n b_n \label{th161:series1}
  \end{gather}
  \begin{enumerate}
    \item $a_n \in \mathbb{R}^n, b_n \in \mathbb{C}, \ n = 1, 2, \dots$
    \item $\{a_n\}, \{a_n\} \downarrow 0 \ (\{a_n\} \uparrow 0)$
    \item $\{B_n\}$ --- последовательность частичных сумм ряда $\sum b_n$ ограничена
  \end{enumerate}
  Тогда ряд \eqref{th161:series1} сходится.
\end{theorem}

\begin{proof}
  $\exists B > 0, \ |B_n| \leq B \ \forall n \Rightarrow \forall m \geq n \geq 2:
  |b_n + \dots + b_m| = |B_m - B_{n - 1}| \leq 2B$ \\
  Возьмем $\varepsilon > 0$. По скольку $a_n \to 0$, то
  $\exists N = N(\varepsilon): \forall n > N(\varepsilon)$ имеем
  $|a_n| < \frac{\varepsilon}{6B}$.\\
  Поэтому, $\forall n > N(\varepsilon)$ и $\forall m \geq n$ получим:
  $$|a_n b_n + \dots + a_m b_m| \leq 2B(|a_n| + 2|a_m|) <
  2B\left(\frac{\varepsilon}{6B} + 2 \frac{\varepsilon}{6B} \right) = \varepsilon$$
  Ряд \eqref{th161:series1} удовлетворяет Критерию Коши сходимости рядов.
\end{proof}
\begin{remark}
  Признак Лейбница - это частный случай признака Дирихле. \\
\end{remark}

\begin{theorem}[признак Абеля]
  \label{th162}
  Если последовательность действительных чисел $a_n$ монотонна и ограничена,
  ряд $\sum\limits_{n = 1}^{\infty} b_n, b_n \in \mathbb{C}$ сходится, то ряд
  \eqref{th161:series1} также сходится.
\end{theorem}

\begin{proof}
  $a_n = a + \alpha_n, \{\alpha_n\}$ --- монотонно стремящаяся к нулю
  последовательность. \\
  Поэтому
  $$\sum a_n b_n = \sum (a + \alpha_n)b_n = a \sum b_n + \sum \alpha_n b_n, $$
  где $a \sum b_n$ сходится по условию, а $\sum \alpha_n b_n$ сходится по
  признаку Дирихле. \\
  $\{B_n\}$ --- последовательность частичных сумм $\sum b_n$ ограничена,
  $\{\alpha_n\}$ --- монотонно стремящаяся к нулю последовательность.
\end{proof}

\section{Безусловно и условно сходящиеся ряды}
\begin{definition}
  Пусть $\{k_n\}, n = 1, 2, \dots$ --- последовательность, в которой каждое
  натуральное число встречается только один раз. $\{k_n\}$ --- однозначное
  отображение $a_n^* = a_{k_n}, \ (n~=~1, 2, \dots)$. \\
  Будем говорить, что ряд $\sum a_n^*$ является перестановкой ряда $\sum a_n$.
\end{definition}
\begin{definition}
  Говорят, что $\sum a_n$ сходится безусловно, если каждая перестановка
  сходится.
\end{definition}
\begin{theorem}
  \label{th171}
  Ряд $\sum a_n, (a_n \in \mathbb{C})$ сходится безусловно тогда и только тогда,
  когда он сходится абсолютно.
\end{theorem}

\begin{proof}
  Достаточность. \\
  Если ряд $\sum a_n$ сходится абсолютно, то все его перестановки сходятся к
  одному и тому же числу --- сумме исходного ряда. \\
  Пусть $\sum a_n^*$ --- перестановка ряда $\sum a_n$. $S_n^*$ --- ее частичная
  сумма. \\
  По Коши: $\forall \varepsilon > 0 \exists N : m \geq n > N$
  \begin{gather}
    |a_n| + \dots + |a_m| < \varepsilon \label{th171:uneq1}
  \end{gather}
  Выберем $p$ так, чтобы все натуральные числа $1, 2, \dots, N$ содержались в
  множестве $k_1, k_2, \dots, k_p$ (смотри определение), тогда при $n > p$
  $a_1, \dots, a_N$ в разности $S_n - S_n^*$ уничтожаются, так что
  $|S_n - S_n^*| < \varepsilon$ в силу \eqref{th171:uneq1}. \\
  Значит $\{S_n^*\}$ сходится к тому же пределу, что и $\{S_n\}$.
\end{proof}

\begin{definition}
  Сходящийся, но не абсолютно сходящийся ряд называется условно сходящимся.
\end{definition}
Из теоремы \eqref{th171} (из необходимости условия) $\Rightarrow$ Теорема \eqref{th172}
\begin{theorem}
  \label{th172}
  Условно сходящийся ряд не может сходится безусловно, то есть у него всегда
  существует расходящаяся перестановка.
\end{theorem}
\begin{proof}
  Без Доказательства.
\end{proof}

\begin{theorem}[Римана]
  Если ряд с действительными членами условно сходится, то каким бы не было
  действительное число $S$, существует перестановка ряда такая, что ее сумма
  равна $S$
\end{theorem}

\begin{proof}
  Без Доказательства.
\end{proof}

  \chapter{Функциональные последовательности и ряды}
  \section{Поточечная сходимость}
Пусть на некотором множестве $\mathbb{E}$ задана последовательность
комплексно значимых функций $f_n, n = 1, 2, \dots \ , \ (f_n \in \mathbb{C})$.
Элементы $x \in \mathbb{E}$ будем называть точками.
\begin{definition}
  $\{f_n\}$ называется ограниченной на $\mathbb{E}$, если
  $\exists M > 0: \forall n \in \mathbb{N}, \forall x \in \mathbb{E}$ выполняется
  $$|f_n(x)| \leq M$$
\end{definition}

\begin{definition}
  $\{f_n\}$ называется сходящейся поточечно на множестве $\mathbb{E}$, если
  при любом фиксированном $x \in \mathbb{E}$, числовая последовательность
  $\{f_n(x)\}$ сходится. Если последовательность сходится на $\mathbb{E}$,
  то $f(x) := \lim\limits_{n \to \infty} f_n(x), x \in \mathbb{E}$ называется
  пределом последовательности. Пусть
  $\{U_n(x)\}_{n = 1}^\infty, x \in \mathbb{E}, \ (U_n \in \mathbb{C})$
  --- последовательность числовых функций.
\end{definition}

\begin{definition}
  Множество числовых рядов
  \begin{gather}
    \sum\limits_{n = 1}^{\infty} U_n(x) \label{def213:series1}
  \end{gather}
  в каждой из которых точка $x$ фиксированная называется рядом на множестве
  $\mathbb{E}$, а функция $U_n(x)$ --- его член. \\
  $S_n(x) = \sum\limits_{k = 1}^{n} U_k(x), x \in \mathbb{E}$ называется
  $n$-ой частичной суммой ряда \eqref{def213:series1}. \\
  $\sum\limits_{k = n + 1}^{\infty} U_k(x)$ - его $n$-ым остатком.
\end{definition}

\begin{definition}
  Ряд \eqref{def213:series1} называется сходящимся поточечно на множестве
  $\mathbb{E}$, если последовательность $\{S_n(x)\}$ сходится поточечно на
  $\mathbb{E}$. При этом $\lim\limits_{n \to \infty} S_n(x) = S(x), x \in \mathbb{E}$
  называется суммой ряда \eqref{def213:series1}.
  $$S(x) = \sum\limits_{n = 1}^{\infty} U_n(x).$$
\end{definition}

\begin{definition}
  Если ряд \eqref{def213:series1} при любом $x \in \mathbb{E}$ сходится абсолютно,
  то он называется абсолютно сходящимся на множестве $\mathbb{E}$.
\end{definition}

\begin{remark}
  Беззаботная перестановка членов ряда может привести к ошибке.
\end{remark}

\section{Равномерная сходимость}
\begin{definition}
  \label{def221}
  Говорят, что функциональная последовательность $\{f_n\}_{n=1}^\infty$
  сходится равномерно на $\mathbb{E}$, если $\forall \varepsilon > 0 \
  \exists N \in \mathbb{N}: \forall n > N, \forall x \in \mathbb{E}$
  имеем
  $$|f_n(x) - f(x)| < \varepsilon$$
  Ясно, что каждая равномерно сходящаяся последовательность, сходится поточечно.
\end{definition}

\begin{comment}
  Обозначение равномерной сходимости:
  $f_n \stackrel{\mathrm{\mathbb{E}}}{\rightrightarrows} f$
\end{comment}

\begin{theorem}[Критерий Коши равномерной сходимости последовательностей]
  \label{th221}
  Для того, чтобы $\{f_n\}$ равномерно сходилась на
  $\mathbb{E} \Longleftrightarrow$
  $\forall \varepsilon > 0 \ \exists N: n, m > N, \forall x \in \mathbb{E}:$
  \begin{gather}
    |f_n(x) - f_m(x)| < \varepsilon \label{th221:uneq1}
  \end{gather}
\end{theorem}

\begin{proof}
  \hfill
  \begin{itemize}
    \item Необходимость: \\
      $f_n \stackrel{\mathrm{\mathbb{E}}}{\rightrightarrows} f$, тогда
      $\forall \varepsilon > 0, \ \exists N \in \mathbb{N}: \forall n > N,
      \forall x \in \mathbb{E} \ |f_n(x) - f(x)| < \frac{\varepsilon}{2}$. \\
      $|f_n(x) - f_m(x)| \leq |f_n(x) - f(x)| + |f(x) - f_m(x)| <
      \frac{\varepsilon}{2} + \frac{\varepsilon}{2} = \varepsilon$,
      $(\forall n, m > N, \forall x \in \mathbb{E}).$
    \item Достаточность: \\
      Пусть выполняется условие Коши, тогда $\{f_n(x)\}$, удовлетворяет
      критерию Коши сходимости числовых последовательностей и следовательно
      сходящегося числового предела, который обозначим $f(x)$. \\
      Тогда перейдя к пределу при $m \to \infty$ получим $\forall n > N,
      \forall x \in \mathbb{E}: |f_n(x) - f(x)| < \varepsilon$. \\
  \end{itemize}
\end{proof}
Иногда полезен критерий, следующий из определения \eqref{def221}

\begin{theorem}
  Пусть $\lim\limits_{n \to \infty} f_n(x) = f(x), \forall x \in \mathbb{E}$. \\
  Положим $r_n = \sup|f_n(x) - f(x)|, x \in \mathbb{E}$ --- равномерное уклонение. \\
  Тогда $f_n \stackrel{\mathrm{\mathbb{E}}}{\rightrightarrows} f
  \Longleftrightarrow r_n \to 0, \ n \to \infty$. (Переформулировка определения).
\end{theorem}

\begin{proof}
  Без доказательства. \\
\end{proof}

\begin{example}
  $f_n(x) = x^n, \mathbb{E} = [0, 1)$ \\
  $\lim\limits_{n \to \infty} f_n(x) = 0, \forall \in \mathbb{E},
  r_n = \sup\limits_{x \in [0, 1)} |x^n - 0| = 1 \not \to 0, n \to \infty.$ \\
  $\{x^n\}$ не является равномерно сходящейся на $\mathbb{E}$.
\end{example}

\begin{example}
  $f_n(x) = x^n - x^{n+1}, \mathbb{E} = [0, 1]$. \\
  $f_n(x) \to 0, \forall x \in \mathbb{E}, \  f_n'(x) = nx^{n-1} - (n + 1)x^n = 0$. \\
  $x_n = \frac{n}{n+1}, \ f_n(x_n) = x_n^n(1 - x_n) < \frac{1}{n + 1}$. \\
  $r_n < \frac{1}{n + 1}$. \\
\end{example}

\begin{definition}
  \label{def222}
  \begin{gather}
    \sum\limits_{n = 1}^{\infty} U_n(x), \ x \in \mathbb{E} \label{def222:series1}
  \end{gather}
  называется равномерно сходящейся, если на множестве $\mathbb{E}$ равномерно
  сходится последовательность частичных сумм. \\
\end{definition}

Пусть $S_k(x)$ --- частичные $k$-ые суммы ряда \eqref{def222:series1},
$$m \geq n: U_n(x) + \dots + U_m(x) = S_m(x) - S_n(x)$$
тогда из теоремы \eqref{th221} (критерий Коши равномерной сходимости последовательности)
$\Rightarrow$ Теорема \eqref{th223} (критерий Коши равномерной сходимости ряда).

\begin{theorem}[Критерий Коши равномерной сходимости ряда]
  \label{th223}
  Для того, чтобы ряд \eqref{def222:series1} равномерно сходился на множестве
  $\mathbb{E} \Longleftrightarrow \forall \varepsilon > 0 \ \exists N \in \mathbb{N},
  \forall n,m > N,
  \forall x \in \mathbb{E}: $
  \begin{gather}
    |U_n(x) + \dots + U_m(x)| < \varepsilon \label{th223:uneq1}
  \end{gather}
\end{theorem}

\begin{proof}
  Без доказательства.
\end{proof}

\begin{consequence}[Необходимый признак равномерной сходимости]
  У равномерно сходящегося ряда общий член равномерно стремится к нулю.
\end{consequence}

\begin{theorem}[Признак Вейерштрасса]
  Пусть $\{U_n\}$ --- последовательнсоть функций, определенных на $\mathbb{E}$
  и пусть $|U_n(x)| \leq a_n, \forall x \in \mathbb{E}, \forall n \in \mathbb{N}.$
  Тогда если $\sum a_n < \infty$ сходится, то следовательно $\sum U_n(x)$ сходится
  равномерно на $\mathbb{E}$.
\end{theorem}
\begin{proof}
  Если $\sum a_n$ сходится, то $\forall \varepsilon > 0 |
  \sum\limits_{k = n}^{m} U_k(x)| \leq
  \sum\limits_{k = n}^{m} a_k < \varepsilon$, при любом $x \in \mathbb{E}$,
  если только $m$ и $n$ достаточно велики, теорема \eqref{th111} (критерий Коши
  сходимости числового ряда). Равномерная сходимость нашего ряда вытекает из
  теоремы \eqref{th223}.
\end{proof}
\begin{remark}
  $\sum a_n$ называется мажорирующим рядом $\sum U_n(x)$.
\end{remark}

\begin{remark}
  ПРОВЕРИТЬ!!! \\
  Условие признака Вейерштрасса не являются необходимыми для равномерной
  сходимости ряда.
\end{remark}

\section{Признаки равномерной сходимости рядов Дирихле и Абеля}
\begin{theorem}
  Пусть дан ряд
  \begin{gather}
    \sum\limits_{n = 1}^{\infty} a_n(x) b_n(x), \ x \in \mathbb{E} \label{th231:series1}
  \end{gather}
  такой что:
  \begin{enumerate}
    \item $a_n(x) \in \mathbb{R}, \ b_n(x) \in \mathbb{C}, \ n = 1, 2, \dots$
    \item $a_n(x) \stackrel{\mathrm{\mathbb{E}}}{\rightrightarrows} 0$
      (Равномерная сходимость к нулю), $\{a_n(x)\}$ - монотонна.
    \item Последовательность частичных сумм $\{B_n(x)\} ряда \ \sum b_n(x)$ ограничена на множестве $\mathbb{E}$.
  \end{enumerate}
  Тогда ряд \eqref{th231:series1} равномерно сходится на множестве $\mathbb{E}$.
\end{theorem}

\begin{proof}
  В силу условия 3, $\exists B > 0: |B_n(x)| \leq B, \ \forall x \in \mathbb{E}, \
  \forall n \in \mathbb{N}$. \\
  $\forall x \in \mathbb{E}, m \geq n \geq 2: |b_n(x) + \dots + b_m(x)| =
  |B_m(x) - B_{n-1}(x)| \leq 2B$. \\
  $\forall \varepsilon > 0$ из условия 2 $\Rightarrow \exists N = N(\varepsilon):
  n > N(\varepsilon), \forall x \in \mathbb{E}$ выполняется неравенство:
  $$0 \leq |a_n(x)| < \frac{\varepsilon}{6B}.$$
  Примениев лемму Абеля \eqref{th151:cons}, получим:
  \begin{gather*}
    |a_n(x) b_n(x) + \dots + a_m(x) b_m(x)| \leq 2B(|a_n(x) + 2a_m(x)|) < \varepsilon \\ 
    \forall x \in \mathbb{E},
    m \geq n \geq N(\varepsilon)
  \end{gather*}
  В силу критерия Коши \eqref{th223}, ряд \eqref{th231:series1} сходится равномерно.
\end{proof}

\begin{theorem}[Признак Абеля]
  \label{th232}
  \begin{gather}
    \sum\limits_{n = 1}^{\infty} a_n(x) b_n(x) \label{th232:series1}
  \end{gather}
  \begin{enumerate}
    \item Если $a_n(x) \in \mathbb{R}, b_n(x) \in \mathbb{C}, n = 1, 2, \dots,
      x \in \mathbb{E}$.
    \item $\{a_n(x)\}$ ограничена на множестве $\mathbb{E}$ и монотонна
      $\forall x \in \mathbb{E}$.
    \item Ряд $\sum b_n(x)$ равномерно сходится на $\mathbb{E}$.
  \end{enumerate}
  Тогда ряд \eqref{th232:series1} равномерно сходится.
\end{theorem}

\begin{proof}
  Доказательство легко провести так, как была доказана теорема \eqref{th161}.
\end{proof}

\begin{example}
  $\sum\limits_{n = 1}^{\infty} \frac{c_n}{n^x}$ ---
  ряд Дирихле. \\
  Если этот ряд сходится в точке $x_0$, то он сходится равномерно
  $\forall x \in \mathbb{E}, \ \mathbb{E} = [x_0, +\infty)$. \\
  Можно воспользоваться Признаком Абеля:
  $$a_n(x) = \frac{1}{n^{x-x_0}}, \ b_n = \frac{c_n}{n^{x_0}}$$
\end{example}

\begin{exercise}
  Рассмотреть и доказать абсолютную сходимость при $x > x_0 + 1$
\end{exercise}

\section{Равномерная сходимость и непрерывность}

\begin{theorem}
  Пусть $f_n \stackrel{\mathrm{\mathbb{E}}}{\rightrightarrows} f, \ x_0$ ---
  предельная точка множества $\mathbb{E}$ и пусть
  $\lim\limits_{x \to x_0} f_n(x) = A_n, \ (n = 1, 2, \dots).$
  Тогда $\{A_n\}$ сходится и
  \begin{gather}
    \lim\limits_{x \to x_0} f(x) = \lim\limits_{n \to \infty} A_n \label{th241:lim1}
  \end{gather}
  Иными словами, 2 предельных перехода в данном случае коммутируют. \\
  $$\lim\limits_{x \to x_0} \lim\limits_{n \to \infty} f_n(x) =
  \lim\limits_{n \to \infty} \lim\limits_{x \to x_0} f_n(x)$$
\end{theorem}

\begin{proof}
  Пусть $\varepsilon > 0.$ В силу равномерной сходимости последовательности
  $\{f_n\} \ \exists N : n > N, m > N, x \in \mathbb{E},$
  \begin{gather}
    \left|f_n(x) - f_m(x)\right| < \varepsilon \label{th241:uneq1}
  \end{gather}
  Переходя в неравенстве \eqref{th241:uneq1} к приделу при $x \to x_0$ получим
  \begin{gather}
    \left|A_n - A_m \right| < \varepsilon, \ (n, m > N) \label{th241:uneq2}
  \end{gather}
  Поэтому $\{A_n\}$ --- последовательность для которой выполняется признак
  Коши сходимости последовательности $\Rightarrow$ она сходится. \\
  Обазначим ее предел $A$ \\
  \begin{gather}
    \left|f(x) - A\right| \leq \left|f(x) - f_n(x)\right| +
    \left|f_n(x) - A_n\right| + \left|A_n - A\right| \label{th241:uneq3}
  \end{gather}
  Выберем $n:$
  \begin{gather}
    |f(x) - f_n(x)| < \frac{\varepsilon}{3}, \ \forall x \in \mathbb{E}
    \label{th241:uneq4}
  \end{gather}
  Это возможно в силу равномерной сходимости.
  \begin{gather}
    |A_n - A| < \frac{\varepsilon}{3} \label{th241:uneq5}
  \end{gather}
  Затем, для этого $n$ подберем такую окрестность
  $U(x_0): x \in U(x_0), x \not = x_0,$ следовательно:
  \begin{gather}
    |f_n(x) - A_n| < \frac{\varepsilon}{3} \label{th241:uneq6}
  \end{gather}
  Из неравенств \eqref{th241:uneq3} --- \eqref{th241:uneq6} получим
  $$|f(x) - A| < \varepsilon, \ \forall x \in U(x_0), \ x \not = x_0$$
  Это равносильно равенству \eqref{th241:lim1}
\end{proof}

\begin{theorem}
  \label{th242}
  Последовательность функций, непрерывных в точке $x \in \mathbb{E}
  f_n \stackrel{\mathrm{\mathbb{E}}}{\rightrightarrows} f,$ то функция
  $f$ непрерывна в точке $x_0$.
\end{theorem}

\begin{proof}
  Без доказательства.
\end{proof}

\begin{remark}
  Обратное не верно, то есть последовательность непрерывных функций может
  неравномерно сходиться.
\end{remark}

Из теоремы \eqref{th242} и определения \eqref{def222} $\Rightarrow$ теорема \eqref{th243}

\begin{theorem}
  \label{th243}
  Если функции $U_n(x), \ (n = 1, 2, \dots), \ x \in \mathbb{E}$ непрерывны
  в точке $x_0 \in \mathbb{E}$ и ряд $\sum\limits_{n = 1}^{\infty} U_n(x)$
  равномерно сходится на $\mathbb{E}$, то его сумма $f(x)$ также непрерывна
  в точкке $x_0$.
\end{theorem}

\section{Равномерная сходимость и интегрирование}
\begin{theorem}
  \label{th251}
  Пусть $f_n$ --- последовательность действительных, значимых, интегрируемых на
  отрезке $[a, b]$ функций. Тогда функция $f$ также интегрируема на $[a,b]$ и
  \begin{gather}
    \int\limits_a^b f(x) dx = \lim\limits_{n \to \infty} \int\limits_a^b f_n(x) dx
    \label{th251:int1}
  \end{gather}
  Существование предела заранее не предполагается.
\end{theorem}

\begin{proof}
  $\forall \varepsilon > 0, \ \exists n:$
  \begin{gather}
    |f_n(x) - f(x)| < \varepsilon, \ x \in [a,b] \label{th251:uneq1}
  \end{gather}
  Зафиксируем $n$ и выберем разбиение
  $[a,b], \ \triangle_1, \dots, \triangle_S$ так, чтобы выполнялось неравенство
  \begin{gather}
    \sum_i \omega(f_n, \triangle_i) |\triangle_i| < \varepsilon \label{th251:sum1}
  \end{gather}
  \begin{comment}
    $\omega(f, E) = \sup - \inf$ --- колебание функции.
  \end{comment}
  Функции $f_n$ интегрируемы на $[a,b]$. По скольку
  $\omega(f, \triangle_i) \leq \omega(f_n, \triangle_i) + 2 \varepsilon,  \ (i = 1, \dots, S)$
  (смотри \eqref{th251:uneq1}).
  $$\sum_i \omega(f, \triangle_i)|\triangle_i| \leq \varepsilon + 2\varepsilon (b - a)$$
  Отсюда следует, что $f \in \mathbb{R} [a,b]$. Для доказательства \eqref{th251:int1}
  выберем $n > N: $
  \begin{gather*}
    |f_n(x) - f(x)| < \varepsilon, \ (a \leq x \leq b), \ n > N \\
    \left|\int\limits_a^b f(x) dx - \int\limits_a^b f_n(x) dx\right| \leq
    \int\limits_a^b |f(x) - f_n(x)| dx < \varepsilon(b - a)
  \end{gather*}
  Отсюда вытекает \eqref{th251:int1}. \\
\end{proof}

\begin{theorem}
  $U_n \in R[a,b]$ (Интегрируема). Если
  \begin{gather}
    f(x) = \sum\limits_{n = 1}^{\infty} U_n(x), \ (a \leq x \leq b)
    \label{th252:series1}
  \end{gather}
  При чем ряд \eqref{th252:series1} сходится на $[a,b]$, тогда
  $$\int\limits_a^b f(x) dx = \sum\limits_{n = 1}^{\infty} \int\limits_a^b
  f(x) dx$$
  Иными словами ряд \eqref{th252:series1} можно интегрировать частями.
\end{theorem}

\begin{proof}
  Без доказательства.
\end{proof}

\begin{remark}
  При нарушении равномерности ряд, состоящий из интегрируемых функций может
  иметь интегрируемую сумму.
\end{remark}

\section{Равномерная сходимость и дифференцирование}
$f_n(x) = \frac{\sin nx}{\sqrt{n}}, x \in \mathbb{R}$ показывает, что из
равномерной сходимости последовательности функций не следует даже поточечная
сходимость последовательностей функций производных. \\
То есть нужны более сильные предположения, чтобы заключать, что $f_n' \to f_n$,
при $f_n \to f$.

\begin{theorem}
  \label{th261}
  Пусть $f_n(x) \to f(x), \ x \in [a,b], \ n \to \infty, \ f_n \in C[a,b], \
  (n = 1, 2, \dots)$. \\
  Если $\{f_n'(x)\}$ сходится равномерно на $[a,b]$, то $f_n(x)$ дифференцируема
  и $$f'(x) = \lim\limits_{n \to \infty} f_n'(x)$$
\end{theorem}

\begin{proof}
  Обозначим через $f^*$ предел последовательности $f_n'$. Ввиду теоремы \eqref{th242}
  $f^*$ непрерывна на $[a,b]$. \\
  Применим теорему \eqref{th251} к последовательости $\{f_n\}$ на промежутке
  $[a, x],$ где $x \in [a,b]$
  $$\int\limits_a^x f^*(t) dt = \lim \int\limits_a^x f'(t) dt =
  \lim\limits_{n \to \infty} (f_n(x) - f_n(a)) = f(x) - f(a)$$
  Так как интеграл слева ввиду непрерывности функции $f^*$ имеет производную
  равную $f',$ то ту же производную имеет и $f(x)$.
  $$f'(x) = f^*(x) = \lim\limits_{n \to \infty} f'(x), x \in [a,b]$$
\end{proof}

Перефразируем теорему \eqref{th261} с точки зрения рядов: \\
Пусть сходящийся ряд $\sum\limits_{n = 1}^{\infty} U_n(x) =: f(x), x \in [a,b]$
и пусть $U_n(x) \in C^1[a,b], \ (n = 1, 2, \dots)$.\\
Если ряд $\sum\limits_{n = 1}^{\infty} U_n'(x)$ сходится равномерно на $[a,b]$,
то сумма $f(x)$ дифференцируема, и \\
$f'(x) = \sum\limits_{n = 1}^{\infty} U_n'(x), x \in [a,b]$.


  \chapter{Степенные ряды}
  \section{Радиус сходимости и круг сходимости}
\begin{definition}
  Степенной ряд --- ряд вида
  \begin{gather}
    \sum\limits_{n = 0}^{\infty} a_n(z - z_0)^n, \ z, z_0 \in \mathbb{C}, n = 0, 1, \dots
    \label{def311:series1}
  \end{gather}
  $a_n$ --- коэффициенты ряда. \\
  $\xi = z - z_0, $ тогда
  $\sum\limits_{n = 0}^{\infty} a_n \xi^n$,
  \begin{gather}
    \sum\limits_{n = 0}^{\infty} a_n z^n \label{def311:series2}
  \end{gather}
\end{definition}

\begin{theorem}
  \label{th311}
  Степенной ряд \eqref{def311:series2}, $\alpha = \overline{\lim} \sqrt[n]{|a_n|}$,
  \begin{gather}
    R = \frac{1}{\alpha} \label{def311:eq1}
  \end{gather}
  ($\alpha = 0 \Longleftrightarrow R = \infty, \
  \alpha = +\infty \Longleftrightarrow, R = 0$), тогда ряд \eqref{def311:series2}
  абсолютно сходится, если $|z| < R,$ и рассходится, если $|z| > R$.
\end{theorem}

\begin{proof}
  Положим $C_n = a_n z^n$. По критерию Коши заключаем, что сумма
  $\sum C_n$ сходится при $\overline{\lim\limits_{n \to \infty}} \sqrt[n]{|C_n|}
  = |z|\cdot \overline{\lim\limits_{n \to \infty}} \sqrt[n]{|a_n|} =
  \frac{|z|}{R} < 1,$ то есть $|z| < R;$ и рассходится, если $|z| > R$.
\end{proof}

\begin{definition}
  Число R называется радиусом сходимости ряда \eqref{def311:series2}. \\
  $|z| < R, z \in \mathbb{C}$ называется кругом сходимости ряда
  \eqref{def311:series2}.
\end{definition}

\begin{remark}
  О сходимсоти на границе окружности $|z| = R$ ничего не говорится в теореме
  \eqref{th311}, так как возможны все варианты.
\end{remark}

\begin{theorem}
  \label{th312}
  Если R --- радиус сходимости $(R > 0)$ ряда \eqref{def311:series2}, то на любом
  круге $|z| < r, $ где $r$ --- фиксированно, и $r < R$. \\
  Таким образом этот ряд сходится абсолютно и равномерно.
\end{theorem}

\begin{proof}
  $z = r, \sum\limits_{n = 0}^{\infty} |a_n| r^n$ сходится, а так как для любой
  точки $z$ круга $|z| \leq r$ выполняется неравенство:
  \begin{gather*}
    |a_n z^n| \leq |a_n| r^n, \ \forall n
  \end{gather*}
  то по признаку Вейерштрассе на этом круге ряд \eqref{def311:series2} сходится
  равномерно.
\end{proof}

\begin{consequence}
  Степеной ряд непрерывный в каждой точке своего круга $|z| < R$ сходится.
\end{consequence}

\begin{theorem}[2-ая т. Абеля]
  \label{th313}
  Если R --- радиус сходимости, $\sum\limits_{n = 0}^{\infty} a_n z^n$ и этот
  ряд сходится при $|z| = R,$ то он сходится на отрезке $[0, R]$ равномерно.
\end{theorem}

\begin{proof}
  Пусть $0 \leq x \leq R$, представим ряд $\sum\limits_{n = 0}^{\infty} a_n x^n
  = \sum\limits_{n = 0}^{\infty} a_n R^n\left(\frac{x}{R}\right)^n$. По скольку
  члены ряда $\sum a_n R^n$ не зависит от $x$, то его сходимость означает его
  равномерную сходимость. $\{(\frac{x}{R})^n\}$ ограничена на отрезке $[0, R]$
  и монотонна в каждой точке. \\
  Поэтому в силу признака Абеля равномерной сходимости рядов \eqref{th232} ряд
  \eqref{def311:series2} равномерно сходится на отрезке $[0, R]$.
\end{proof}

\begin{lemma}
  \label{ch3:lemma1}
  Радиусы сходимости $R, R_1, R_2$ соответственно рядов
  $\sum\limits_{n = 0}^{\infty} a_n z^n, \sum\limits_{n = 0}^{\infty}
  \frac{a_n}{n + 1} z^{n+1}, \sum\limits_{n = 0}^{\infty} n a_n z^{n - 1}$ равны:
  $R = R_1 = R_2$.
\end{lemma}

\begin{proof}
  Действительно, так как $\lim\limits_{n \to \infty} \sqrt[n]{\frac{1}{n + 1}}=
  \lim\limits_{n \to \infty} \sqrt[n]{n} = 1$, то \\
  $\overline{\lim\limits_{n \to \infty}} \sqrt[n]{|a_n|} =
  \overline{\lim\limits_{n \to \infty}} \sqrt[n]{\frac{a_n}{n + 1}}=
  \overline{\lim\limits_{n \to \infty}} \sqrt[n]{|n a_n|}$
\end{proof}

\begin{example}
  $\sum a_n (z - z_0)^n$.
  Областью сходимости такого ряда является круг $|z - z_0| < R$, с точностью до
  граничных точек.
\end{example}

\section{Степенные ряды в действительной области. Общие свойства.}
В параграфах 3.2 - 3.4 будем рассматривать
\begin{gather}
  \sum\limits_{n = 0}^{\infty} a_n (x - x_0)^n,
  \label{ch3:lim1}
\end{gather}
где $a_n, x, x_0$ --- действительные числа. \\
Если $R$ --- радиус сходимости ряда ряда \eqref{ch3:lim1}, то очевидно ряд
\eqref{ch3:lim1} сходится, если $|x| < R$ и расходится, если $|x| > R$. \\
Число $R$ --- по-прежнему называется радиусом сходимости ряда \eqref{ch3:lim1},
а интервал \\ $(x_0 - R, x_0 + R)$ --- его интервал сходимости.

\begin{theorem}
  \label{th321}
  Если $R$ --- радиус сходимости ряда
  \begin{gather}
    f(x) = \sum\limits_{n = 0}^{\infty} a_n (x - x_0)^n,
    \label{ch3:lim2}
  \end{gather}
  где $R > 0$, то:
  \begin{enumerate}
    \item функция $f$ имеет в интервале $(x_0 - R, x_0 + R)$ производные всех
      порядков, они называются почленным диффиринциалом ряда \eqref{ch3:lim2}:
      \begin{gather}
        f^{(m)}(x) = \sum\limits_{n = m}^{\infty} n(n-1)\dots(n - m + 1)
        a_n(x - x_0)^{n-m}, \ m = 1, 2, \dots
        \label{ch3:lim3}
      \end{gather}
    \item $\forall x \in (x_0 - R, x_0 + R)$
      \begin{gather}
        \int\limits_{x_0}^x f(t) dt = \sum\limits_{n = 0}^{\infty} a_n
        \frac{(x - x_0)^{n+1}}{n+1}
        \label{ch3:lim4}
      \end{gather}
    \item \eqref{ch3:lim2} - \eqref{ch3:lim4} имеют одинаковые радиусы сходимости
      $R$.
  \end{enumerate}
\end{theorem}

\begin{proof}
  В силу леммы \eqref{ch3:lemma1} ряды \eqref{ch3:lim3}, \eqref{ch3:lim4} имеют
  тот же радиус сходимости, что и ряд \eqref{ch3:lim2}. Всякий ряд с $R >
  0$ сходится на отрезке $[x_0 - r, x_0 + r]$, \\
  $0 < r < R \ $(теорема~\eqref{th312}). \\
  Поэтому утверждения 1 и 2 непосредственно следуют из общих теорем о
  сходимости рядов (\eqref{th151:cons} и \eqref{th162}).
\end{proof}

\begin{theorem}
  Если функция $f$ раскладывается в некоторой окружности $x_0$, то она
  раскладывается в степенной ряд.
  \begin{gather*}
    f(x) = \sum\limits_{n = 0}^{\infty} a_n (x - x_0)^n
  \end{gather*}
  \begin{gather}
    a_n = \frac{f^{(n)}(x_0)}{n!}, \ n = 0, 1, \dots \label{ch3:lim5} \\
    f(x) = \sum\limits_{n = 0}^{\infty} \frac{f^{(n)}(x_0)}{n!} (x - x_0)^n
    \label{ch3:lim6}.
  \end{gather}
\end{theorem}

\begin{consequence}
  \label{th322:cons}
  Если в некоторой окружности точки функция раскладывается в степенной ряд, то
  это разложение единственно.
\end{consequence}

\begin{proof}
  Продифференцировав $m$ раз равенство \eqref{ch3:lim2}, получим (в силу \eqref{ch3:lim3}):

  \begin{gather*}
    f^{(m)}(x) = m(m-1)\dots 2 \cdot 1 \cdot a_m +
    (m+1)m \dots a_{m-1} (x-x_0) + (m+2)(m+1) \dots 3 \cdot a_{m-2} (x -
    x_0)^2\dots
  \end{gather*}
  Положим $x = x_0$, тогда получаем:
  \begin{gather*}
    f^{(m)}(x_0) = m! \ a_m, \ m = 0, 1, \dots
  \end{gather*}
\end{proof}

\section{Ряд Тейлора. Разложение функции в степенные ряды.}
\begin{definition}
  Пусть $f$ определена в некоторой окрестности точки $x_0$ и имеет в этой точке
  производные всех порядков, тогда ряд
  \begin{gather}
    \sum\limits_{n = 0}^{\infty} \frac{f^{(n)}(x_0)}{n!} (x-x_0)^n
    \label{def331:series1}
  \end{gather}
  Называется рядом Тейлора функции $f$ в точке $x_0$.
\end{definition}

Следующий пример показывает, что функция, бесконечно дифференцируемая в одной
точке может быть не равна разложению по Тейлору в окрестности этой точки.

\begin{example}
  \begin{gather*}
    f(x) =
    \begin{cases}
      e^{-1/x^2}, \ x \not = 0 \\
      0, \ x = 0
    \end{cases}
  \end{gather*}
  $f^{(n)}(0) = 0, \ n = 0, 1, \dots$ \\
  Отсюда следует, что все члены ряда Тейлора \eqref{ch3:lim2} в точке $x_0 = 0$,
  и не совпадают с функцией $f(x)$ в никакой окрестности точки $x_0$.
\end{example}

\begin{approval}
  \label{app331}
  Пусть функция $f(x)$ определена в некоторой окрестности $(x_0 - h, x_0 + h)$.
  \begin{gather}
    S_n(x) = \sum\limits_{k = 0}^{n} \frac{f^{(k)}(x_0)}{k!} (x - x_0)^k
    \label{app331:sum1} \\
    r_n(x) = f(x) - S_n(x) \label{app331:term1}
  \end{gather}
  Тогда, для того, чтобы функция $f(x)$ на интервале $(x_0 - h, x_0 + h)$ равна
  сумме своего ряда~\eqref{def311:series1}, то есть:
  \begin{gather}
    (S_n(x) \to f(x), \ n \to \infty) \Longleftrightarrow \lim\limits_{n \to
    \infty} r_n(x) = 0, \ \forall x \in (x_0 - h, x_0 + h)
    \label{app331:lim1}
  \end{gather}
\end{approval}

\begin{theorem}
  \label{th331}
  Пусть функция $f$ и все ее производные ограничены в совокупности на интервале
  $(x_0 - h, x_0 + h)$, то есть существует такая $M = const, M > 0$: \\$\forall
  x \in (x_0 - h, x_0 + h), \ n = 0, 1, \dots, $ выполняется неравенство:
  \begin{gather}
    |f^{(n)}(x)| \leq M \label{th331:uneq1}
  \end{gather}
  Тогда на интервале $(x_0 - h, x_0 + h)$ функция $f$ раскладывается в ряд
  Тейлора:
  \begin{gather}
    f(x) = \sum\limits_{n = 0}^{\infty} \frac{f^{(n)}(x_0)}{n!} (x - x_0)^n,
    \label{th331:series1}
  \end{gather}
  где $|x - x_0| < h$.
\end{theorem}

\begin{proof}
  \begin{gather}
    \forall a : \lim\limits_{n \to \infty} \frac{a^n}{n!} = 0 \label{th331:lim1}
  \end{gather}
  По формуле Тейлора с остаточным членом в форме Лагранжа, для $x \in (x_0 - h,
  x_0 + h)$, для $\forall M$ имеем:
  \begin{gather*}
    f(x) = S_n(x) + r_n(x),
  \end{gather*}
  где $r_n(x) = \frac{f^{(n+1)}(\xi)}{(n+1)!}(x - x_0)^{n+1}$, где $\xi = x_0 +
  \theta(x - x_0)$, где $0 < \theta < 1$. \\
  Используя \eqref{th331:uneq1} получим:
  \begin{gather*}
    |r_n(x)| = \frac{|f^{(n+1)}(\xi) (x - x_0)^{n+1}|}{(n+1)!} \leq
    \frac{M|x-x_0|^{n+1}}{(n+1)!}, \ \forall x \in (x_0 - h, x_0 + h).
  \end{gather*}
  Остюда из \eqref{th331:lim1} следует \eqref{app331:lim1}. Согласно
  утверждению \eqref{app331} теорема доказана.
\end{proof}

\section{Разложение основных элементарных в ряд Тейлора.}
\begin{itemize}
  \item Разложение в ряд функции $e^x, \cos x, \sin x$. \\
    Использую теорему \eqref{th331}, получаем:
    \begin{gather*}
      f^{(n)}(x) = e^x, \ \sin(x + \frac{\pi}{2}n), \ \cos(x + \frac{\pi}{2}n),
      \ n = 0, 1, \dots,
    \end{gather*}
    Так что $|f^{(n)}(x)| \leq e^h, \ f(x) = e^x, |x| \leq h$ \\
    $|f^{(n)}(x)| \leq 1, \ f(x) = \sin x, \cos x, \forall x \in \mathbb{R}$ \\
    Так как коэффициенты Тейлора для этих функций известны, то мы можем
    записать разложение при любом $x$:
    \begin{gather}
      e^x = \sum\limits_{n = 0}^{\infty} \frac{x^n}{n!} \label{ch34:series1} \\
      \sin x = \sum\limits_{n = 0}^{\infty} (-1)^n \frac{x^{2n+1}}{(2n+1)!}
      \label{ch34:series2} \\
      \cos x = \sum\limits_{n = 0}^{\infty} (-1)^n \frac{x^{2n}}{(2n)!}
      \label{ch34:series3}
    \end{gather}
  \item Разложение в ряд функции $\sh x, \ch x$. \\
    Заменив в \eqref{ch34:series1} $x$ на $-x$ получим
    \begin{gather}
      e^{-x} = \sum\limits_{n = 0}^{\infty} \frac{(-1)^n x^n}{n!}
      \label{ch34:series4}
    \end{gather}
    Отсюда из \eqref{ch34:series1} получаем:
    \begin{gather}
      \sh x = \frac{1}{2} \left(e^x - e^{-x}\right) = \sum\limits_{n =
      0}^{\infty} \frac{x^{2n+1}}{(2n + 1)!} \label{ch34:series5} \\
      \ch x = \frac{1}{2} \left(e^x + e^{-x}\right) = \sum\limits_{n =
      0}^{\infty} \frac{x^{2n}}{(2n)!} \label{ch34:series6}
    \end{gather}
    В правых частях этих формул разложения степенных функций в ряды единственно
    в силу теоремы \eqref{th322:cons}.
  \item Разложение в ряд функции $\ln(1 + x)$. \\
    Рассмотрим:
    \begin{gather}
      \frac{1}{1 + t} = 1 - t + t^2 - t^3 + \dots + (-1)^n t^n + \dots, \ |t| < 1
      \label{ch34:series7}
    \end{gather}
    Интегрирую его почленно по теореме \eqref{th321} от $0$ до $x \in (-1, 1)$
    получим:
    \begin{gather*}
      \int\limits_0^x \frac{dt}{1+t} = \ln(1+x) = x - \frac{x^2}{2} +
      \frac{x^3}{3} - \dots, \\
      \ln(1+x) = \sum\limits_{n = 1}^{\infty} \frac{(-1)^{n+1} x^n}{n}, \
      \forall x \in (-1, 1)
      \label{ch34:series8}
    \end{gather*}
    Ряд правой части равенства \eqref{ch34:series8} сходится по признаку Лейбница
    $\Rightarrow$ согласно теореме Абеля \eqref{th313}, разложение
    \eqref{ch34:series8} имеет место в промежутке $(-1, 1]$
  \item Разложение в ряд $(1 + x)^\alpha, \alpha \not = 0, 1, \dots$
    Формула Тейлора для этой функции имеет вид:
    \begin{gather}
      (1+x)^\alpha = 1 + \alpha x + \frac{\alpha(\alpha - 1)}{2!}x^2 + \dots +
      \frac{\alpha(\alpha - 1)\dots(\alpha - n + 1)}{n!} x^n + r_n(x)
      \label{ch34:series9}
    \end{gather}
    Соответствующий степенной ряд называют
    \begin{gather}
      1 + \sum\limits_{n = 1}^{\infty} \frac{\alpha(\alpha - 1)\dots(\alpha
      -n + 1)}{n!} x^n \label{ch34:series10}
    \end{gather}
    биномиальным рядом. \\
    $R = \lim\limits_{n \to \infty} \left|\frac{a_n}{a_{n+1}}\right| =
    \lim\limits_{n \to \infty} \left|\frac{n+1}{\alpha - n}\right| = 1$,
    в силу утверждения, $r_n(x) \to 0$.

    \begin{remark}
      Поведение ряда \eqref{ch34:series10} в точках $\pm 1$, характерезуется
      следующей таблицой:
      \begin{table}[H]
        \caption{таблица, характеризующая ряд \eqref{ch34:series10}}
        \begin{center}
          \begin{tabular}{|c|c|l|}
            \hline
            & $\alpha > 0$ & абсолютно сходится \\
            $x = 1$ & $-1 < \alpha < 0$ & условно сходится \\
            & $\alpha \leq -1$ & расходится\\
            \hline
            $x = -1$ & $ \alpha > 0$ & абсолютно сходится \\
            & $\alpha < 0$ & рассходится \\
            \hline
          \end{tabular}
        \end{center}
      \end{table}
      Согласно второй теореме Абеля \eqref{th313} всякий раз, когда ряд
      \eqref{ch34:series10} сходится при $x = \pm 1$, его сумма равна $(1 +
      x)^\alpha$.
    \end{remark}
  \item Разложение в ряд $\arctg x$ \\
    Рассмотрим ряд:
    \begin{gather}
      \frac{1}{1 + t^2} = 1 - t^2 + t^4 - t^6 + \dots + (-1)^n t^{2n} + \dots, \ |t| < 1
    \end{gather}
    Интегрирую его почленно по теореме \eqref{th321} от $0$ до $x \in (-1, 1)$
    получим:
    \begin{gather*}
      \int\limits_0^x \frac{dt}{1+t^2} = \arctg x = x - \frac{x^3}{3} +
      \frac{x^5}{5} - \dots + (-1)^n \frac{x^{2n+1}}{(2n+1)} + \dots
    \end{gather*}
    Ряд правой части равенства \eqref{ch34:series8} сходится по признаку Лейбница
    $\Rightarrow$ согласно теореме Абеля \eqref{th313}, разложение
    \eqref{ch34:series8} имеет место на отрезке $(-1, 1)$. \\
    В частности, при $x = 1$, получим:
    \begin{gather*}
      1 - \frac{1}{3} + \frac{1}{5} - \frac{1}{7} + \dots = \frac{\pi}{4}
    \end{gather*}
  \item Разложение в ряд $\arcsin x$ \\
    Рассмотрим ряд:
    \begin{gather}
      \frac{1}{\sqrt{1 - t^2}} = 1 + \sum\limits_{n = 1}^{\infty}
      \frac{(2n-1)!!}{(2n)!!}t^{2n}, \ |t| < 1
    \end{gather}
    Интегрирую его почленно по теореме \eqref{th321} от $0$ до $x \in (-1, 1)$
    получим:
    \begin{gather*}
      \int\limits_0^x \frac{dt}{\sqrt{1-t^2}} = \arcsin x = x + \sum\limits_{n
      = 1}^{\infty} \frac{(2n-1)!!}{(2n)!!} \frac{x^{2n+1}}{2n+1}, \ |x| \leq 1
    \end{gather*}
    Справедливость этого разложения при $x = \pm 1$ устанавливается с помощью
    второй теоремы Абеля \eqref{th313}.
\end{itemize}

\section{Формулы Эйлера}
Ряды разложения \eqref{ch34:series1} - \eqref{ch34:series3} функций $e^x, \sin x,
\cos x$ сходятся всюду в комплексной плоскости $\mathbb{C}$. По этой причине
естественны следующие определения($e^z, \sin z, \cos z, z \in \mathbb{C}$):
\begin{gather}
  e^z = \exp(z) = \sum\limits_{n = 0}^{\infty} \frac{z^n}{n!}
  \label{ch35:series1} \\
  \sin z = \sum\limits_{n = 0}^{\infty} \frac{(-1)^n z^{2n+1}}{(2n+1)!}
  \label{ch35:series2} \\
  \cos z = \sum\limits_{n = 0}^{\infty} \frac{(-1)^n z^{2n}}{(2n)!}
  \label{ch35:series3}
\end{gather}
Заменив $z$ сначала на $iz$, а затем на $-iz$ получим:
\begin{gather}\label{ch35:series4}
  e^{iz} = \sum\limits_{n = 0}^{\infty} \frac{i^n z^n}{n!} \\
  e^{-iz} = \sum\limits_{n = 0}^{\infty} \frac{(-1)^n i^n z^n}{n!}
\end{gather}

Заметим, что $i^{2k} = (-1)^k, i^{2k + 1} = (-1)^ki, k = 0, 1, \dots$. \\
\begin{gather*}
  \frac{1}{2} (e^{iz} + e^{-iz}) = \sum\limits_{n = 0}^{\infty} \frac{(-1)^n
  z^{2n}}{(2n)!} \\
  \frac{1}{2i} (e^{iz} - e^{-iz}) = \sum\limits_{n = 0}^{\infty}
  \frac{(-1)^n z^{2n+1}}{(2n+1)!}
\end{gather*}
Сравнив эти формулы с \eqref{ch35:series2}, \eqref{ch35:series3} заключаем, что
\begin{gather}
  \cos z = \frac{e^{iz} + e^{-iz}}{2}
  \label{ch35:eq1} \\
  \sin z = \frac{e^{iz} - e^{-iz}}{2i}
  \label{ch35:eq2}
\end{gather}

Из этих формул следует формула:
\begin{gather}
  \cos z + i \sin z = e^{iz} \label{ch35:eq3}
\end{gather}
Формулы \eqref{ch35:eq1}, \eqref{ch35:eq2} и \eqref{ch35:eq3} называются
формулами Эйлера. \\
Если в формуле \eqref{ch35:eq3} $z = \varphi, \ \varphi \in
\mathbb{R}$, то
\begin{gather*}
  \cos \varphi + i \sin \varphi = e^{i\varphi}
\end{gather*}
Поэтому $z \in \mathbb{C}, |z| = r, \ z = r(\cos \varphi + i \sin
\varphi)$
\begin{gather*}
  z = r e^{i\varphi}
\end{gather*}

\begin{definition}
  $W(x) = U(x) + iV(x), \ x \in \mathbb{R}, V(x) \in \mathbb{R}$ \\
  Положим $\frac{dW}{dx} = U'(x) + iV'(x)$, тогда
  \begin{gather*}
    \int\limits_a^b W(x) dx = \int\limits_a^b U(x) dx + i \int\limits_a^b V(x)
    dx
  \end{gather*}
\end{definition}
$e^{i\pi} = -1$

  %\chapter{Радомский консультация}
  %\input{arad}
  \chapter{Ряды Фурье}
  \section{Ортогональные системы} \label{ch41}
В параграфах \eqref{ch41} - \eqref{ch43} $\mathbb{X}$ --- линейное
бесконечномерное пространство(действительное или комплексное, со скалярным
произведением).
$$X(\cdot \ , \ \cdot), \ \|x\| = \sqrt{(x, x)}.$$
$\mathbb{K}$ --- некоторое счетное или конечное множество.
\begin{definition}
  \label{def411}
  Система векторов $\{x_k: k \in \mathbb{K}\}, \ x \in \mathbb{X}$ ---
  ортогональная система(ОС). \\
  $(x_i, x_j) = 0, \ \forall i, j \in \mathbb{K}, i \not = j$ (и система не
  нулевая). Если $(x_i, x_i) = 1$, то система называется ортонормированной.
\end{definition}

\begin{theorem}
  Ортогональная система векторов линейно независима, то есть линейно не
  зависима каждая ее конечная подсистема.
\end{theorem}
\begin{proof}
  Определение линейной независимости:
  \begin{gather*}
    \alpha_1 x_1 + \dots + \alpha_i x_i +
    \dots = 0 \Longleftrightarrow \alpha_i = 0, \ \forall i
  \end{gather*}
  Скалярно умножим все члены на $x_i$, тогда получим:
  \begin{gather*}
    \label{th411:eq1}
    \alpha_1 (x_1, x_i) + \dots + \alpha_i (x_i, x_i) + \dots = (0, x_i) \\
    \label{th411:eq2}
    \alpha_i (x_i, x_i) = 0 \\
    \label{th411:eq3}
    \alpha_i = 0
  \end{gather*}
  Равенство \eqref{th411:eq2} следует из определения \eqref{def411}, равенство
  \eqref{th411:eq3} следует из того, что $(x_i, x_i)~\not=~0$ (так как система
  не нулевая).
\end{proof}

\section{Коэффициенты Фурье}
\begin{definition}
  Пусть $\{e_k: \ k \in \mathbb{K}\}$ --- ОНС в $\mathbb{X}$, $\{(x, e_k)\}, x
  \in \mathbb{X}$ называется коэффициентами Фурье элемента $x$ в ОНС $e_k$.
\end{definition}
\begin{lemma}
  \label{lemma421}
  Если система векторов $e_1, \dots, e_n$ пространства $\mathbb{X}$ --- ОН, то
  $\forall x \in \mathbb{X}$ вектор $h = x - x_e$, где
  \begin{gather}
    x_e = \sum\limits_{k = 1}^{n} (x, e_k) e_k
    \label{lemma1:eq1}
  \end{gather}
  ортоганален подпространству $\mathbb{L} = \langle e_1, \dots, e_n \rangle$ (натянотому
  на векторы $e_1, \dots, e_n$)
\end{lemma}

\begin{proof}
  Достаточно проверить, что скалярное произведение \\ $(h, e_j) = 0, \ \forall j =
  1, \dots, n$
  \begin{gather*}
    (h, e_j) = (x, e_j) - \sum\limits_{k = 1}^{n} (x, e_k) (e_k, e_j) =
    (x, e_j) - (x, e_j) = 0
  \end{gather*}
\end{proof}

\begin{lemma}[теорема Пифагора]
  \label{lemma422}
  Если векторы $x_1, \dots, x_n$ попарно ортогональны и $x = x_1 + \dots +
  x_n$, то $\|x\|^2 = \|x_1\|^2 + \dots + \|x_n\|^2$
\end{lemma}

\begin{proof}
  $(x, x) = (\sum\limits_{i = 1}^{n} x_i, \sum\limits_{i = 1}^{n} x_i) =
  \sum\limits_{i, j = 1}^{n} (x_i, x_j) = \sum\limits_{i = 1}^{n} (x_i, x_i)$
\end{proof}

\begin{theorem}[экстремальное свойство коэффициентов Фурье]
  \label{th423}
  Если $e_1, \dots, e_n$ --- ОНС пространства $\mathbb{X}$, то $\forall x \in
  \mathbb{X}$ и $\forall y = \alpha_1 e_1 + \dots + \alpha_n e_n$ имеет место
  неравенство:
  \begin{gather*}
    \|x - \sum\limits_{k = 1}^{n} (x, e_k) e_k \| \leq \|x - \sum\limits_{k =
    1}^{n} \alpha_k e_k \|,
  \end{gather*}
  в котором равенство возможно при условии: $\alpha_k = (x, e_k) \ \forall k = 1,
  \dots, n$.
\end{theorem}

\begin{proof}
  Представим $x - y$ в виде $x - y = (x_e - y) + h$, где $x_e, h$ определены в
  лемме~\eqref{lemma421}. \\
  По лемме~\eqref{lemma421} $h \perp (x_e - y) \in \mathbb{L}$. По теореме
  Пифагора (лемма~\ref{lemma422}):
  \begin{gather*}
    \|x - y\|^2 = \|x_e - y\|^2 + \|h\|^2 = \|x_e - y\|^2 + \|x - x_e\|^2 \geq
    \|x - x_e\|^2
  \end{gather*}
  равенство возможно, когда коэффициенты $\alpha_k$ совпадают с коэффициентами
  Фурье.
\end{proof}

\begin{remark}
  Теорема \eqref{th423} показывает, что вектор $x_e$ является наилучшей в
  смысле нормы пространства $\mathbb{X}$, аппроксимацией вектора $x$
  подпространства $\mathbb{L} = \langle e_1, \dots, e_n \rangle$, так что
  наименьшее уклонением вектора $x$ от $\mathbb{L}$ равно $\|x - x_e\|$.
\end{remark}

\begin{theorem}[неравенство Бесселя]
  \label{th424}
  Если $\{e_1, \dots, e_n\}$ --- ОНС в $\mathbb{X}$, то $\forall x \in
  \mathbb{X}$ справедливо неравенство:
  \begin{gather}
    \sum\limits_{k = 1}^{n} |(x, e_k)|^2 \leq \|x\|^2
    \label{th424:uneq1}
  \end{gather}
  Если $\{e_k: k \in \mathbb{K}\}$ --- ОНС, то $\forall x \in \mathbb{X}$
  \begin{gather}
    \sum |(x, e_k)|^2 \leq \|x\|^2
    \label{th424:uneq2}
  \end{gather}
\end{theorem}

\begin{proof}
  По лемме \eqref{lemma421}
  \begin{gather*}
    x = \sum\limits_{k = 1}^{n} (x, e_k)e_k + h,
  \end{gather*}
  при чем система векторов $e_1, \dots, e_n, h$ --- ортогональна в $\mathbb{X}$
  \\
  по теореме Пифагора получаем:
  \begin{gather}
    \|x\|^2 = \sum\limits_{k = 1}^{n} |(x, e_k)|^2 + \|h\|^2
    \label{th424:eq1}
  \end{gather}
  остюда следует \eqref{th424:uneq1}, (так как это имеет место для любой
  конечной системы векторов), отсюда следует \eqref{th424:uneq2}.
\end{proof}

\begin{remark}
  Из \eqref{th424:eq1} следует формула наименьшего отклонения:
  \begin{gather}
    \|x - x_e\|^2 \equiv \|x - \sum\limits_{k = 1}^{n} (x, e_k) e_k \|^2 =
    \|x\|^2 - \sum\limits_{k = 1}^{n} |(x, e_k)|^2
    \label{th424:eq3}
  \end{gather}
\end{remark}

Переформулируем понятие коэффициентов Фурье для произвольной ОС(не обязательно
нормированной) $\{f_k\}$. \\
Для этого по этой системе построим ОНС. \\
$\{e_k = \frac{f_k}{\|f_k\|}\}$, используем ортгональное разложение:
\begin{gather*}
  x = x_e +
  h, \ x = \sum\limits_{k = 1}^{n} (x, e_k) e_k + h = \sum\limits_{k =
  1}^{n} \frac{(x, f_k)}{\|f_k\|^2} f_k + h
\end{gather*}

\begin{definition}
  \label{def422}
  $\{e_k = \frac{(x, f_k)}{\|f_k\|^2} \}$ --- называется коэффициентами Фурье
  вектора $x$ в ОС~$\{f_k\}$. \\
  Заменим в неравенстве \eqref{th424:uneq2}, $e_k$ на $\frac{f_k}{\|f_k\|}$
  получим неравенство Бесселя для произвольной ОС.
  \begin{gather}
    \sum\limits_{k \in \mathbb{K}} \frac{|(x, f_k)|^2}{\|f_k\|^2} \leq \|x\|^2,
    \ \{f_k, k \in \mathbb{K}\}
    \label{def422:uneq1}
  \end{gather}
  Или, в других обозначения:
  \begin{gather*}
    \sum\limits_{k \in \mathbb{K}} |C_k|^2 \|f_k\|^2 \leq \|x\|^2
  \end{gather*}
\end{definition}

\begin{example}
  В пространстве $\mathbb{X} = \mathcal{R}_2([-\pi, \pi], \mathbb{C})$. \\
  Рассмотрим ОС $\{e^{ikt}: k \in \mathbb{Z}\}$. \\
  В соответствии с определением \eqref{def422} коэффициенты Фурье $C_k$ функции
  $f$ в системе $\{e_{ik}\}$ выражаются формулами:
  \begin{gather}
    C_k = \frac{1}{2\pi} \int\limits_{-\pi}^\pi f(t) e^{-ikt} dt
    \label{ex421:coef1}
  \end{gather}
  из неравенства Бесселя \eqref{def422:uneq1} $\forall f \in
  \mathcal{R}_2([-\pi,\pi], \mathbb{C})$
  \begin{gather}
    \sum\limits_{k = -\infty}^{\infty} |C_k|^2 \leq \frac{1}{2\pi}
    \int\limits_{-\pi}^\pi |f(t)|^2 dt
    \label{ex421:uneq1}
  \end{gather}
\end{example}

\begin{example}
  Аналогично находим коэффициенты Фурье. \\
  $\{\frac{1}{2} a_0, a_k, b_k : k \in \mathbb{N}\}$ функции $f \in
  \mathcal{R}([-\pi, \pi], \mathbb{C})$ в ОС $\{1, \cos kx, \sin kx : k \in
  \mathbb{N}\}$:
  \begin{gather}
    a_k = \frac{1}{\pi} \int\limits_{-\pi}^\pi f(t) \cos kt dt, \ k = 0, 1,
    \dots \label{ex422:coef1} \\
    b_k = \frac{1}{\pi} \int\limits_{-\pi}^\pi f(t) \sin kt dt, \ k = 1, 2,
    \dots \label{ex422:coef2}
  \end{gather}
  по неравенству Бесселя все принимает вид:
  \begin{gather}
    \frac{|a_0|^2}{2} + \sum\limits_{k = 1}^{\infty} (|a_k|^2 + |b_k|^2) \leq
    \frac{1}{\pi} \int\limits_{-\pi}^\pi |f(t)|^2 dt \label{ex422:uneq1}
  \end{gather}
\end{example}

\begin{remark}
  Сравнивая равенства \eqref{ex422:coef1}, \eqref{ex422:coef2} и
  \eqref{ex421:coef1} с учетом формулы Эйлера получаем:
  \begin{gather}
    C_k =
    \begin{cases}
      \frac{1}{2} (a_k - ib_k), \ k \geq 0 \\
      \frac{1}{2} (a_k + ib_{-k}), \ k < 0
    \end{cases}
    \label{rm423:cases1}
  \end{gather}
\end{remark}

\section{Ряд Фурье}
\label{ch43}
\begin{definition}
  Если $\{f_1, \dots, f_k, \dots\}$ --- ОС в $\mathbb{X}$, а $x \in
  \mathbb{X}$, то можно составить ряд:
  \begin{gather*}
    x \sim \sum\limits_{k = 1}^{\infty} C_k f_k,
  \end{gather*}
  где $C_k = \frac{(x, f_k)}{\|f_k\|^2}$. \\
  Этот ряд называется рядом Фурье вектора $x$ по ОС $\{f_k\}$. \\
  Ряд Фурье по ОНС $\{e_k\}$ имеет вид:
  \begin{gather*}
    x \sim \sum\limits_{k = 1}^{\infty} (x, e_k) e_k
  \end{gather*}
\end{definition}

\begin{definition}
  Говорят, что ряд $\sum\limits_{k = 1}^{\infty} y_k, y_k \in \mathbb{X}$
  сходится в $\mathbb{X}$ к вектору $x \in \mathbb{X}$ (сходится по
  норме($\|\cdot\|$) пространства $\mathbb{X}$), если
  \begin{gather*}
    \lim\limits_{n \to \infty} \|x - \sum\limits_{k = 1}^{n} y_k\| = 0
  \end{gather*}
  При этом пишем $x \overset{\mathbb{X}}= \sum\limits_{k = 1}^{\infty}
  y_k$ по норме пространства $\mathbb{X}$.
\end{definition}

\begin{theorem}
  \label{th431}
  $\{e_k: k \in \mathbb{N}\}$ --- ОНС в $\mathbb{X}, x \in \mathbb{X}$, где \\
  $x \overset{\mathbb{X}} = \sum\limits_{k = 1}^{\infty} (x, e_k) e_k
  \Longleftrightarrow$ когда $\|x\|^2 = \sum\limits_{k = 1}^{\infty} |(x,
  e_k|^2$. \\
  Это равенство называется равенством Парсеваля и представляет собой обобщение
  теоремы Пифагора на случай бесконечномерного пространства.
\end{theorem}

\begin{definition}
  Система $\{x_k: k \in \mathbb{K} \}$ векторов в пространстве $\mathbb{X}$
  называется полной в множестве $\mathbb{E} \subset \mathbb{X}$, если любой
  вектор $x \in E$ можно сколь угодно точно в смысле нормы пространства
  $\mathbb{X}$ приблизить к конечной линейной комбинации векторов системы.
\end{definition}

\begin{theorem}
  \label{th432}
  Пусть $\{e_1, \dots, e_n, \dots\}$ --- ОНС в $\mathbb{X}$, тогда следующее
  условие эквивалентны:
  \begin{enumerate}
    \item $\{e_k\}$ полна в множестве $\mathbb{E} \subset \mathbb{X}$.
    \item $\forall x \in \mathbb{E} \subset \mathbb{X}$ имеет место разложение
      (в ряд Фурье) $x \overset{\mathbb{X}} = \sum\limits_{k = 1}^{\infty} (x,
      e_k) e_k$
    \item $\forall x \in \mathbb{E} \subset \mathbb{X}$ имеет место равенство
      Парсеваля: $\|x\|^2 = \sum\limits_{k = 1}^{\infty} |(x, e_k)|^2$.
  \end{enumerate}
\end{theorem}

\begin{proof}
  Из $1 \Rightarrow 2$ в силу экстремального свойства коэффициентов Фурье. \\
  Из $2 \Rightarrow 3$ по теореме \eqref{th431} \\
  Из $3 \Rightarrow 1$, по скольку по формуле уклонений
  \begin{gather*}
    \|x - \sum\limits_{k = 1}^{n} (x, e_k)\|^2 = \|x\|^2 - \sum\limits_{k =
    1}^{n} |(x, e_k)|^2 \to 0, \text{при} \ n \to \infty
  \end{gather*}
\end{proof}

\section{Тригонометрический ряд Фурье}
$\mathbb{X} = \mathbb{R}_2([-\pi, \pi], \mathbb{C}), e_k = \{e^{ikx} : k \in
\mathbb{Z}\}, f \in \mathbb{R}([-\pi, \pi], \mathbb{C})$
\begin{gather}
  C_k(f) = C_k = \frac{1}{2} \int\limits_{-\pi}^\pi f(x) e^{-ikx} dx
  \label{ch44:eq1}
\end{gather}
Сопоставим функцию
\begin{gather}
  f(x) \sim \sum\limits_{k = -\infty}^{+\infty} C_k e^{ikx}
  \label{ch44:eq2}
\end{gather}

\begin{definition}
  Если нам дан тригонометрический ряд Фурье в комплексной записи, то его
  $n$-ая частная сумма равна:
  \begin{gather}
    S_n(x) = S_n(f, x) = \sum\limits_{k = -n}^{n} C_k e^{ikx}
    \label{def441:eq1}
  \end{gather}
\end{definition}

\begin{definition}
  Ряд Фурье функции $f \in \mathbb{R}([-\pi, \pi], \mathbb{C})$ по системе
  $\{1, \cos kx, \sin kx: k \in \mathbb{N}\}$ называется тригонометрическим
  рядом Фурье и записывается следующим образом:
  \begin{gather}
    f(x) \sim \frac{a_0}{2} + \sum\limits_{k = 1}^{\infty} (a_k \cos kx + b_k
    \sin kx)
    \label{def442:sim1} \\
    a_k = \frac{1}{\pi} \int\limits_{-\pi}^\pi f(x) \cos kx dx, \ k = 0, 1,
    \dots
    \label{def442:coef1} \\
    b_k = \frac{1}{\pi} \int\limits_{-\pi}^\pi f(x) \sin kx dx, \ k = 1, 2,
    \dots
    \label{def442:coef2}
  \end{gather}
  Если функция $f$ --- действительная, то $a_k, b_k \in \mathbb{R}$ и
  $\underline{C_k} = \overline{C_k} \ (k = 0, 1, \dots)$.
\end{definition}

\begin{definition}
  Тригонометрические многочлены $D_n(x) = \sum\limits_{k = -n}^{n} e^{ikx},
  K_n(x) = \frac{1}{n + 1} \sum\limits_{m = 0}^{n} D_m(x)$ называются
  соответственно ядром Дирихле и ядром Фейера.
  \begin{gather*}
    a_1, \dots, a_n \to a, \ \text{при} \ n \to \infty \\
    \frac{a_1 + \dots + a_n}{n} \to a, \ \text{при} \ n \to \infty
  \end{gather*}
  а обратное не верно.
\end{definition}

\begin{gather}
  D_n(x) = \sum\limits_{k = -n}^{n} e^{ikx}
  \label{ch44:kernels} \\
  K_n(x) = \frac{1}{n + 1} \sum\limits_{m = 0}^{n} D_m(x)
\end{gather}

\begin{theorem}
  При $n = 0, 1, \dots$ имеем:
  \begin{gather}
    D_n(x) = \frac{\sin((n + \frac{1}{2})x)}{\sin \frac{x}{2}}
    \label{th441:eq1} \\
    K_n(x) = \frac{1}{n + 1} \frac{1 - \cos ((n + 1)x)}{1 - \cos x}
    \label{th441:eq2} \\
    \frac{1}{\pi} \int\limits_{-\pi}^\pi D_n(x) dx = \frac{1}{2\pi}
    \int\limits_{-\pi}^\pi K_n(x) dx = 1
    \label{th441:eq3}
  \end{gather}
  кроме того,
  \begin{gather}
    K_n(x) \geq 0, K_n(x) \leq \frac{2}{(n+1)(1 - \cos \delta)},
    \text{где} \ 0 < \delta \leq |x| \leq \pi
    \label{th441:uneq1}
  \end{gather}
\end{theorem}

\begin{proof}
  Согласно \eqref{ch44:kernels}
  \begin{gather}
    D_n(x) = \frac{e^{i(n+1)x} - e^{-inx}}{e^{ix} - 1}
    \label{th441:eq4}
  \end{gather}
  чтобы получить \eqref{th441:eq1} домножим здесь числитель и знаменатель на
  $e^{-ix/2}$. \\
  Подставим \eqref{th441:eq4}, в определение ядра $K_n(x)$, получим:
  \begin{gather*}
    K_n(x) = \frac{1}{n + 1} \cdot \frac{e^{-ix} - 1}{(e^{-ix} - 1)(e^{ix} - 1)}
    \sum\limits_{m = 0}^{n} (e^{i(m+1)x} - e^{-imx}) = \\ \frac{1}{n + 1} \cdot
    \frac{1}{2 - (e^{ix} + e^{-ix})} \sum\limits_{m = 0}^{n} [(e^{imx} +
    e^{-imx}) - (e^{i(m+1)x} + e^{-i(m+1)x})] = \\ \frac{1}{n + 1}
    \cdot \frac{1}{2 - 2
    \cos x} (2 - (e^{i(m+1)x} + e^{-i(m+1)x}))
  \end{gather*}
  откуда следует \eqref{th441:eq2}. \\
  Значит $K_n(x) \geq 0$ и выполняется \eqref{th441:uneq1}. А \eqref{th441:eq3}
  непосредственно следует из \eqref{ch44:kernels}.
\end{proof}

Далее предпологаем, что функция $f$, изначально определенная на $[-\pi, \pi]$,
продолжена на $\mathbb{R}$ как $2\pi$-периодическая функция. \\
Если $f \in C[-\pi, \pi]$, то ее $2\pi$-периодическое продолжение непрерывно на
$\mathbb{R}$

\begin{gather*}
  (f \in C_{2\pi}) \Longleftrightarrow f(-\pi) = f(\pi)
\end{gather*}

\begin{lemma}[интегральное представление частичной суммы ряда Фурье]
  \begin{gather}
    \forall x \in \mathbb{R}: \ S_n(f, x) = \frac{1}{2\pi}
    \int\limits_{-\pi}^\pi f(x - t) D_n(t) dt
    \label{lem441:eq1}
  \end{gather}
\end{lemma}

\begin{proof}
  Пусть $S_n$ --- частичная сумма, тогда:
  \begin{align*}
    S_n(f, x) = \sum\limits_{k = -n}^{n} C_k e^{ikx} &= \sum\limits_{k = -n}^{n}
    \frac{1}{2\pi} \int\limits_{-\pi}^\pi f(u) e^{-iku} du \cdot e^{ikx} = \\
    \frac{1}{2\pi} \int\limits_{-\pi}^\pi f(u) \sum\limits_{k = -n}^{n} e^{ik(x
    - u)} du &=  \frac{1}{2\pi} \int\limits_{-\pi}^\pi f(u) D_n(x - u) du = \\
    \frac{1}{2\pi} \int\limits_{x -\pi}^{x + \pi} f(x - t) D_n(t) dt &=
    \frac{1}{2\pi} \int\limits_{-\pi}^\pi f(x - t) D_n(t) dt
  \end{align*}
  \begin{comment}
    Использовалась замена: $x - u = t$.
  \end{comment}
  Отметим, что последнее равенство выполняется, поскольку в следствии
  периодичности функции безразлично по какому инетрвалу интегрировать, лишь бы
  его длина была равна~$2\pi$.
\end{proof}

\begin{definition}
  Средние арифмитические частичных сумм:
  \begin{gather}
    S_n(f, x), \delta_n(f, x) = \delta_n(x) = \frac{S_0(x) + \dots +
    S_n(x)}{n+1}
    \label{def444:eq1}
  \end{gather}
  Называются полиномами Фейера.
\end{definition}

\begin{theorem}[теорема Фейера]
  Если функция $f \in C_{2\pi}$, то
  \begin{gather}
    \delta_n(x) \stackrel{\mathrm{\mathbb{R}}}{\rightrightarrows} f(x)
    \label{th442:fnef1}
  \end{gather}
\end{theorem}

\begin{proof}
  Согласно формулам \eqref{def444:eq1}, \eqref{lem441:eq1} и
  \eqref{ch44:kernels} имеем: \\
  \begin{gather}
    \delta_n(x) = \frac{1}{2\pi} \int\limits_{-\pi}^\pi f(x - t) K_n(t) dt
  \end{gather}
  поэтому из \eqref{th441:eq3} следует, что:
  \begin{gather}
    f(x) - \delta_n(x) = \frac{1}{2\pi} \int\limits_{-\pi}^\pi (f(x) - f(x -
    t)) K_n(t) dt
    \label{th442:eq1}
  \end{gather}
  $\varepsilon > 0, \ M = \max |f(x)|, \ x \in \mathbb{R}$. Поскольку функция
  $f$ --- равномерно непрерывна, то найдется такое $\delta > 0$, что:
  \begin{gather}
    |x - y| < \delta \Rightarrow |f(x) - f(y)| < \frac{\varepsilon}{2}
    \label{th442:cons1}
  \end{gather}
  Согласно \eqref{th441:uneq1} можно затем выбрать такое $N = N(\varepsilon,
  \delta)$, что
  \begin{gather}
    n > N, \delta \leq |t| \leq \pi \Rightarrow K_n(t) \leq
    \frac{\varepsilon}{4M}
    \label{th442:uneq1}
  \end{gather}
  Из \eqref{th442:cons1} и $K_n(t) \geq 0$ получаем:
  \begin{gather}
    \int\limits_{-\delta}^\delta |f(x) - f(x - t)| |K_n(t)| dt <
    \frac{\varepsilon}{2} \int\limits_{-\pi}^\pi K_n(t) dt = \pi \varepsilon, \
    n = 1, 2, \dots
    \label{th442:uneq2} \\
    \left\{\int\limits_{-\pi}^\delta + \int\limits_{\delta}^\pi \right\} =
    |f(x) - f(x-t)| |K_n(t)| dt \leq \frac{\varepsilon}{4M}
    \int\limits_{-\pi}^\pi 2M dt = \pi \varepsilon, \ \forall n > N
    \label{th442:uneq3}
  \end{gather}
  В силу \eqref{th442:cons1}, \eqref{th442:uneq2} и \eqref{th442:uneq3}
  получаем:
  \begin{gather}
    |f(x) - \delta_n(x)| < \varepsilon, \ \forall x \in \mathbb{R}, \ \forall n
    > N
    \label{th442:uneq4}
  \end{gather}
\end{proof}

\begin{consequence}
  Если две непрерывные $2\pi$-периодические функции $f, g$ имеют один и тот же
  ряд Фурье, то $f(x) = g(x), \ \forall x \in \mathbb{R}$
\end{consequence}

\begin{proof}
  Действительно, если $\delta_n(x)$ --- среднее арифметическое этого ряда, то:
  \begin{gather*}
    \delta_n(x) \to f(x), \ \delta_n \to g(x)
  \end{gather*}
\end{proof}

\begin{consequence}
  Если $f \in C_{2\pi}$ и $\int\limits_{-\pi}^\pi f(x) e^{-inx} dx \equiv 0$,
  то $\forall n \in \mathbb{Z}: \ f(x) \equiv 0$. Таким образом ОС $\{e^{ikx}:
  \ k \in \mathbb{Z}\}$ нельзя дополнить ненулевым элементом.
\end{consequence}

\begin{proof}
  Это вытекает из предыдущего следствия, если положить $g = 0$.
\end{proof}

\begin{consequence}
  Ряд Фурье функции $f \in C_{2\pi}$ либо сходится в каждой точке $x$ к функции
  $f(x)$, либо вовсе расходится в этой точке.
\end{consequence}

\begin{proof}
  НАПИСАТЬ ДОКАЗАТЕЛЬСТВО.
\end{proof}

\begin{remark}
  Ряд Фурье для выражения функции в самом деле может в некоторых точках
  расходится.
\end{remark}

\begin{theorem}[теорема Вейерштрасса]
  \label{th443}
  Если $f \in C_{2\pi}$, то $\forall \varepsilon > 0$ существует такой
  тригонометрический многочлен $T(x): \ \forall x \in \mathbb{R}$:
  \begin{gather*}
    |f(x) - T(x)| < \varepsilon
  \end{gather*}
\end{theorem}

\begin{proof}
  Без доказательства.
\end{proof}

Из теоремы \eqref{th443} следует теорема \eqref{th444}.
\begin{theorem}[теорема Вейерштрасса]
  \label{th444}
  Если $f \in C[a, b]$, то $\forall \varepsilon > 0$ существует такой
  алгеброический многочлен $P(x): \ \forall x \in [a, b]$:
  \begin{gather*}
    |f(x) - P(x)| < \varepsilon
  \end{gather*}
\end{theorem}

\begin{proof}
  Положив $t = \frac{x - a}{b - a} \pi,$ и $x = \frac{b - a}{\pi}t + a$, получим
  функцию $\varphi(t) = f(a + \frac{b - a}{\pi}t)$ на отрезке $[0, \pi]$.
  Продолжим ее в начале четным образом $\varphi(-t) = \varphi(t), \ t \in [\pi,
  0)$. Найдем по теореме \eqref{th443} такой тригонометрический полином
  $T(x): |\varphi(t) - T(t)| < \frac{\varepsilon}{2}, \ \forall t \in
  \mathbb{R}$. Всякий тригонометрический полином раскладывается по Тейлору,
  сходится равномерно на любом конечном интервале. \\
  Пусть $P_n$ --- частичная сумма ряда Тейлора для $T(t)$ такая что: $|T(t) -
  P_n(t)| < \frac{\varepsilon}{2}, \ 0 \leq t \leq \pi$. Тогда $|\varphi_n(t) -
  P_n(t)| < \varepsilon$, при $0 \leq t \leq \pi$. Сделав обратную замену в
  $P_n(t): t = \frac{x - a}{b - a} \pi$, получим многочлен $Q_n(x)$,
  удовлетворяющий условию: $|f(x) - Q_n(x)| < \varepsilon, \ a \leq x \leq b$.
\end{proof}

Получаем еще одно следствие от теоремы Фейера --- полнота тригонометрической
системы функций $C_2[-\pi, \pi]$ и более общо $\mathcal{R}_2[-\pi, \pi]$.

\begin{theorem}[о полноте тригонометрической системы]
  \label{th445}
  Любая функция $f$ из множества $f \in \mathcal{R}[-\pi, \pi]$ может быть сколь угодно
  точно приближена в среднем, то есть по норме:
  \begin{enumerate}
    \item Кусочно-постоянной функции $[-\pi, \pi]$
    \item Непрерывными на отрезке $[-\pi, \pi]$ функциями, принимающие равные
      значения на концах $[-\pi, \pi]$
    \item Тригонометрическими полиномами
  \end{enumerate}
\end{theorem}

\begin{proof}
  Достаточно рассмотреть случай действительно значимых функций.
  \begin{enumerate}
    \item Поскольку $f$ --- интегрируема, то: $\forall \varepsilon > 0 \
      \exists$ разбиение $-\pi = x_0 < x_1 < \dots < x_n = \pi$ отрезка $[-\pi,
      \pi]$, что: $0 \leq \triangle := \int\limits_{-\pi}^\pi f(x) dx -
      \sum\limits_{i = 1}^{n} m_i \triangle x_i < \varepsilon$, где
      \begin{gather*}
        m_i = \inf\{f(x)\}, \ x \in [x_{i-1}, x_i), \ \triangle x_i = x_i - x_{i-1}
      \end{gather*}
      Полагая
        $g(x) =
        \begin{cases}
          m_i, \text{если} \ x \in [x_{i-1}, x_i) \\
          0, \text{если} \ x = \pi \\
        \end{cases} $,
      $M_f = \sup\{|f(x)|\}, |x| \leq \pi$. \\
      Получим:
        $\int\limits_{-\pi}^\pi (f(x) - g(x))^2 dx \leq \int\limits_{-\pi}^\pi
        (|f(x)| + |g(x)|) (|f(x)| - |g(x)|) dx \leq \\
        2M_f \int\limits_{-\pi}^\pi
        (f(x) - g(x)) dx = 2M_f\triangle \leq 2M_f \varepsilon$.
      \item Достаточно уметь приблежать к среднему кусочно постоянной функции.
        Пусть $g$ --- такая функция, с точками разрыва $x_1, \dots, x_n$.
        Удобно присваивать $-\pi = x_1, x_n = \pi$. \\
        Очевидно, какого бы ни было $\varepsilon > 0 \ \exists \delta > 0$, что
        $\delta$-окрестности точек $x_1, \dots, x_n$ не пересекаются и
        $2\delta_n M < \varepsilon$, где $M = \sup \{|g(x)| : |x| \leq \pi\}$
        Заменим функцию $g$ на каждом из отрезков $[-\pi, -\pi + \delta], [x_1
        -\delta, x_1 + \delta], (i = 2, \dots, n-1), [\pi - \delta, \pi]$
        линейной функцией, принимающей на концах этих отрезков соответственно: \\
        $0, g(-\pi + \delta), g(x_i - \delta), g(x_i + \delta), (i = 2, \dots,
        n - 1), g(\pi - \delta), 0$. Получим кусоно линейную, непрерывную на
        кусочном отрезке $[-\pi, \pi]$ функцию $h, h(-\pi) = h(\pi) = 0, \\
        |h(x)| \leq M, \ \forall x \in [-\pi, \pi]$. \\
        Значит
        \begin{gather}
          \int\limits_{-\pi}^\pi (g - h)^2 dx \leq 2M
          \int\limits_{-\pi}^\pi (|g - h|) dx = \\ 2M \sum\limits_{i = 1}^{n}
          \int\limits_{x_i - \delta}^{x_i + \delta} (|g - h|)dx \leq 2M (2M -
          2\delta)n < 4M\varepsilon
        \end{gather}
      \item Осталось показа, что можно приблизить любую функцию класса 2. Но по
        теореме Фейера для любой функции типа $h$, найдется такой
        тригонометрический многочлено, что:
        \begin{gather*}
          \forall \varepsilon > 0, \ T: |h(x) - T(x)| < \varepsilon, \ \forall
          x \in [-\pi, \pi] \\
          \left(\frac{1}{2\pi} \int\limits_{-\pi}^\pi (h(x) - T(x))^2 dx
          \right)^{\frac{1}{2}} < \varepsilon
        \end{gather*}
        Ссылаясь на неравенство треугольника, пространства $\mathcal{R_2}[-\pi, \pi]$
        заключаем, что теорема доказана.
  \end{enumerate}
\end{proof}

Из полноты тригонометрической системы, из теоремы \eqref{th432} (третьего
условия полноты ОС) и форумлы наименьших уклонений следует теорема
\eqref{th446}.

\begin{theorem}
  \label{th446}
  $f \in \mathcal{R}_2([-\pi, \pi], \mathbb{C})$, имеем:
  \begin{enumerate}
    \item $f(x) \overset{\mathcal{R}_2}=\frac{a_0(f)}{2} + \sum\limits_{k = 1}^{\infty}
      a_k(f) \cos kx + b_k(f) \sin kx$, или в комплексной записи: \\
      $f(x) \overset{\mathcal{R}_2}= \sum\limits_{-\infty}^{+\infty} C_k(f) e^{ikx}$, где
      сходимость понимается, как сходимость по норме.
    \item $\frac{1}{\pi} \int\limits_{-\pi}^\pi |f(x)|^2 dx =
      |\frac{a_0(f)}{2}|^2 + \sum\limits_{k = 1}^{\infty} |a_k(f)|^2 +
      |b_k(f)|^2$, или в комплексной записи:\\
      $\frac{1}{2\pi} \int\limits_{-\pi}^\pi |f(x)|^2 dx =
      \sum\limits_{-\infty}^{+\infty} |C_k(f)|^2$.
    \item $\frac{1}{\pi}\int\limits_{-\pi}^\pi |f(x) - S_n(f, x)|^2 dx =
      \sum\limits_{k = n + 1}^{\infty} |a_k(f)|^2 + |b_k(f)|^2 = 4
      \sum\limits_{k = n + 1}^{\infty} |C_k(f)|^2, f \in \mathcal{R}([-\pi, \pi],
      \mathbb{R})$.
  \end{enumerate}
\end{theorem}

\section{Обобщение на неограниченные функции}
\begin{definition}
  Пусть $0 < p < \infty, -\infty \leq a < b \leq + \infty$. Будем писать $f \in
  \mathcal{R}^p[a,b]$, если существует конечное число точек $x_j, j = 0, 1, 2, \dots, n$,
  таких что:
  \begin{enumerate}
    \item $a = x_0 < x_1 < \dots < x_n = b$
    \item Функция $f$ интегрируема по Римману на любом отрезке $f \in
      \mathcal{R}[\alpha, \beta], \ [\alpha, \beta] \subset (x_{j-1}, x_j)$
    \item Интеграл $\int\limits_{x_{j-1}}^{x_j} |f(x)|^p dx, \ j = 1, 2,
      \dots, n$ сходится.
  \end{enumerate}
\end{definition}

\begin{remark}
  Формулы \eqref{ch44:eq1}, \eqref{def442:coef1} и \eqref{def442:coef2},
  определяющие коэффициенты Фурье $C_k(f), a_k(f), b_k(f)$ имеют смысл для
  $\forall f \in \mathcal{R}^1[-\pi, \pi]$, так как тогда $f(x) e^{-ikx}, f(x) \cos kx,
  f(x) \sin kx \in \mathcal{R}^1[-\pi, \pi]$.
\end{remark}

\begin{definition}[неравенство Гельдера]
  \begin{gather*}
    \int\limits_a^b |f(x) g(x)| dx \leq \left(\int\limits_a^b |f(x)|^p dx \right)
    ^{\frac{1}{p}}
    \left(\int\limits_a^b |g(x)|^q dx\right)^{\frac{1}{q}}
  \end{gather*}
  где $q > 1, \ p > 1, \ \frac{1}{p} + \frac{1}{q} = 1$
  \begin{gather*}
    \int\limits_a^b |f(x)| dx \leq (b - a)^{1 - \frac{1}{p}}
    \left(\int\limits_a^b |f(x)|^p dx \right)^{\frac{1}{p}} \\
    \left(\frac{1}{b-a} \int\limits_a^b |f(x)|^{r} dx \right)
    ^{\frac{1}{r}} \leq \left(\frac{1}{b-a} \int\limits_a^b |f(x)|^{h} dx
    \right) ^{\frac{1}{h}}
  \end{gather*}
  если $h > 1, \ 0 < r < h < \infty$.
\end{definition}

\begin{approval}
  теоремы \eqref{th445} и \eqref{th446} остаются в силе, если в них
  пространство $\mathcal{R}_2[-\pi, \pi]$ расширить до линейного пространства
  $\mathcal{R}^2[-\pi, \pi]$:
  \begin{gather*}
    (f, g) = \int f(x) \overline{g(x)} dx
  \end{gather*}
\end{approval}

\begin{proof}
  Смотри теорему \eqref{th445}.
\end{proof}

\begin{lemma}
  \label{lem451}
  Если $f \in \mathcal{R}^p[a,b] (p > 0)$, то $\forall \varepsilon > 0 \exists g \in
  \mathcal{R}[a,b]: \int\limits_a^b |f(x) - g(x)|^p dx < \varepsilon$.
\end{lemma}

\begin{proof}
  Доказать самостоятельно.
\end{proof}

\section{Достаточные условия сходимости тригонометрического ряда Фурье в точке}
\begin{lemma}[Римана]
  Если $f \in \mathcal{R}^1[a,b]$, то
  \begin{gather}
    \int\limits_a^b f(x) e^{i\lambda x} dx \to 0
    \label{lem461:lim1} \\
    \int\limits_a^b f(x) \cos \lambda x dx \to 0
    \label{lem461:lim2} \\
    \int\limits_a^b f(x) \sin \lambda x dx \to 0, \text{при} \ \lambda \to
    \infty, \ \lambda \in \mathbb{R}
    \label{lem461:lim3}
  \end{gather}
\end{lemma}

\begin{proof}
  Будем считать, что функция $f(x)$ --- действительная, так как в случае $f(x)$
  --- комплексная легко сводится к этому, согласно лемме \eqref{lem451} при $p
  = 1$, $\forall \varepsilon > 0, \ \exists g$ (кусочно-постоянная):
  \begin{gather}
    \int\limits_a^b |f(x) - g(x)| dx < \frac{\varepsilon}{2}
    \label{th461:uneq1}
  \end{gather}
  Пусть $g(x) = m, x \in [x_{j-1}, x_j], j = 1, \dots, n$. $(x_0 = a, x_n =
  b)$, тогда:
  \begin{gather*}
    \int\limits_a^b g(x) e^{i\lambda x} dx = \sum\limits_{j = 1}^{n}
    \int\limits_{x_{j-1}}^{x_j} m_j e^{i\lambda x} dx = \left. \frac{1}{ix}
    \sum\limits_{j = 1}^{n} (m_j e^{i \lambda x})
    \right |_{x_{j-1}}^{x_j} \to 0, \ \text{при} \ \lambda \to \infty
  \end{gather*}
  Отсюда из \eqref{th461:uneq1} получим: \\
  \begin{gather*}
    \left| \int\limits_a^b f(x) e^{i\lambda x} dx \right| =
    \left| \int\limits_a^b (f(x) - g(x)) e^{i\lambda x} dx +
    \int\limits_a^b g(x) e^{i\lambda x} dx \right| \leq \\
    \int\limits_a^b |f(x) - g(x)| e^{i\lambda x} dx +
    \left| \int\limits_a^b g(x) e^{i\lambda x} dx \right| \leq
    \frac{\varepsilon}{2} + \frac{\varepsilon}{2} = \varepsilon
  \end{gather*}
  Итак, \eqref{lem461:lim1} доказано. \\
  Отделяя действительные и мнимые части, получаем \eqref{lem461:lim2},
  \eqref{lem461:lim3}
\end{proof}

\begin{definition}
  Говорят, что функция $f$, заданная в проколотой окрестности точки $x$,
  удовлетворяет условиям Дини, если при $x \in \mathbb{R}$ выполняется:
  \begin{enumerate}
    \item В точке $x$ существуют оба предела:
      \begin{gather*}
        f(x-0) = \lim\limits_{t \to +0} f(x - t) \\
        f(x+0) = \lim\limits_{t \to +0} f(x + t)
      \end{gather*}
    \item $\int\limits_0^{\delta} \frac{f(x - t) - f(x - 0)}{t} dt \ $,
      $\int\limits_0^{\delta} \frac{f(x + t) - f(x + 0)}{t} dt$ сходятся
      абсолютно на $[0, \delta], \ \forall \delta > 0$
  \end{enumerate}
\end{definition}

\begin{theorem}
  $f$ --- $2\pi$-периодическая функция $f \in \mathcal{R}^1[-\pi, \pi], \ x \in
  \mathbb{R}$. \\
  \begin{gather}
    \sum\limits_{\infty} C_k(f) e^{ikx} = \frac{f(x-0) + f(x+0)}{2}
    \label{th461:eq1}
  \end{gather}
\end{theorem}

\begin{proof}
  $S_n(f, x) = \frac{1}{2\pi} \int\limits_{-\pi}^\pi f(x-t) D_n(t) dt =
  \frac{1}{2\pi} \int\limits_{-\pi}^0 f(x-t) D_n(t) dt + \frac{1}{2\pi}
  \int\limits_{0}^\pi f(x-t) D_n(t) dt = \frac{1}{2\pi} \int\limits_{0}^\pi
  f(x+t) D_n(t) dt + \frac{1}{2\pi} \int\limits_{0}^\pi f(x-t) D_n(t) dt =
  \frac{1}{2\pi} \int\limits_{0}^\pi (f(x - t) + f(x + t)) D_n(t) dt = \\
  \frac{1}{\pi} \int\limits_0^\pi \frac{f(x-0) + f(x+0)}{2} D_n(t) dt +
  \frac{1}{\pi} \int\limits_0^\pi \left(\frac{f(x-t) - f(x-0)}{2} +
  \frac{f(x+t) - f(x+0)}{2}\right) D_n(t) dt = \\
  \frac{f(x-0) + f(x+0)}{2} + \frac{1}{\pi} \int\limits_0^\pi \left(
  \frac{f(x-t) - f(x-0)}{2\sin \frac{t}{2}} + \frac{f(x+t) - f(x+0)}{2\sin
  \frac{t}{2}} \right) \sin((n+\frac{1}{2}) t) \ dt
  \label{th461:eqs}$ \\
  Поскольку, $2 \sin \frac{t}{2} \sim t, \ t \to 0$, то из условий Дини
  следует, что $g_x(t) \in \mathcal{R}^1[0, \pi]$ абсолютно интегрируема. На основании
  леммы Римана:
  \begin{gather*}
    \int\limits_0^\pi g_x(t) \sin ((n + \frac{1}{2})t) dt \to 0, \text{при} \ n
    \to \infty
  \end{gather*}
  Отсюда из \eqref{th461:eqs} следует \eqref{th461:eq1}
\end{proof}

\begin{consequence}
  Пусть $f$ --- ограниченная функция с периодом $2\pi$, имеющая разрыв первого
  рода и пусть имеет левые и правые производные. Тогда ряд Фурье сходится
  всюду, а его сумма в точке разрыва непрерывна и равна: $f(x) =
  \frac{1}{2}(f(x+0) + f(x-0))$
\end{consequence}

\section{Гладкость функции и скорость убывания коэффициентов Фурье}
\begin{definition}
  Функцию $f$ называют кусочно диффиренцируемой, если существует ее разбиение
  $a = x_0 < x_1 < \dots < x_n = b$, такое что на каждом интервале $(x_{j-1},
  x_j) f$ --- диффиренцируема, а в точках $x_{j-1}, x_j$ существуют конечные
  значения $f(x_{j-1}+0), f(x_j-0), f_+'(x_{j-1}), f_-'(x_j)$, при $\forall j =
  1, \dots, n, \ f' \in Q[a,b]$ \\
  Через $C_{2\pi}^{(k)} (k = 0, 1, \dots)$ обозначим класс $2\pi$-периодических
  (комплексно-значимых) функций, имеющих на $\mathbb{R} \ k$-ую непрерывную
  производную. $f^{(k)}, \ (f^{(k)} \in C_{2\pi} = C_{2\pi}^{(0)})$.
\end{definition}

\begin{lemma}[о диффиренцировании ряда Фурье]
  Если $f$ --- непрерывна, $f \in C_{2\pi}$ и $f$ кусочно диффиренцируема на
  $[-\pi, \pi]$, то $f' \sim \sum\limits_{-\infty}^{+\infty} C_k(f') e^{ikx}$
  может быть получен формальным диффиренцированием ряда Фурье $f \sim
  \sum\limits_{-\infty}^{+\infty} C_k(f) e^{ikx}$ самой функции, то есть:
  \begin{gather}
    C_k(f') = ikC_k(f), \ k \in \mathbb{Z}
    \label{lem471:eq1}
  \end{gather}
\end{lemma}

\begin{proof}
  Интегрированием по частям находим:
  \begin{gather*}
    C_k(f') = \frac{1}{2\pi} \int\limits_{-\pi}^\pi f'(x) e^{-ikx} dx =
    \left. \frac{1}{2\pi} f(x) e^{-ikx} \right|_{-\pi}^\pi + \frac{ik}{2\pi}
    \int\limits_{-\pi}^\pi f(x) e^{-ikx} dx = ikC_k(f)
  \end{gather*}
\end{proof}

\begin{theorem}
  Пусть $f \in C_{2\pi}^{(m-1)}, \ (m \in \mathbb{N}), \ f^{(m-1)}$ ---
  кусочно-диффиренцируема \\ на $[-\pi, \pi]$. Тогда:
  \begin{gather}
    C_k(f^{(m)}) = (ik)^m C_k(f), \ k \in \mathbb{Z}
    \label{th471:eq1} \\
    |C_{\pm k}(f)| = \frac{\gamma_k}{k^m} =
    \bar{\bar{o}}\left(\frac{1}{k^m}\right), \ k
    \to \infty, k \in \mathbb{N}
    \label{th471:eq2}
  \end{gather}
  Причем $\sum\limits_{k = 1}^{\infty} \gamma_k^2 < \infty, \ \gamma_k \downarrow
  0, \ \gamma_k = \bar{\bar{o}}\left(\frac{1}{k}\right)$.
\end{theorem}

\begin{proof}
  Соотношение \eqref{th471:eq1} получается в результате $m$-кратного
  использования \eqref{lem471:eq1}. \\
  $C_k(f^{(m)}) = ik C_k(f^{(m-1)}) = \dots = (ik)^m C_k(f), \ k \in
  \mathbb{Z}$. \\
  Полагаю $\gamma_k = |C_k(f^{(m)})|$, с учетом неравенства Бесселя:
  \begin{gather*}
    \sum\limits_{-\infty}^{+\infty} \gamma_k^2 \leq \frac{1}{2\pi}
  \end{gather*}
  Из \eqref{th471:eq1} получаем \eqref{th471:eq2}.
\end{proof}

\begin{theorem}
  $f \in C_{2\pi}^{(m-1)}, \ m \in \mathbb{N}, \ f^{(m-1)}$ ---
  кусочно-диффиренцируема на $[-\pi, \pi]$, тогда ряд Фурье сходится абсолютно
  и равномерно на $\mathbb{R}$, причем
  \begin{gather}
    \max\limits_{x \in \mathbb{R}} |f(x) - S_n(f,x)| =
    \bar{\bar{o}}\left(\frac{1}{n^{m - 1/2}}\right), \ n \to \infty
    \label{th472:eq1}
  \end{gather}
\end{theorem}

\begin{proof}
  Поскольку $f \in C_{2\pi}$ и удовлетворяет условиям Дини в любой точке, то
  \begin{gather*}
    f(x) = \lim\limits_{n \to \infty} S_n(f, x), \ \forall x \in \mathbb{R}
  \end{gather*}
  Поэтому, с учетом \eqref{th471:eq2} имеем:
  \begin{gather}
    |f(x) - S_n(f, x)| = \left|\sum\limits_{|k| \geq n + 1} C_k(f)
    e^{ikx}\right| \leq \left|\sum\limits_{|k| \geq n + 1}
    |C_k(f)e^{ikx}|\right| = \\
    2 \sum\limits_{k = n + 1}^{\infty} \frac{\gamma_k}{k^m} \leq
    2 \left(\sum\limits_{k = n + 1}^{\infty} \gamma_k^2 \right)^{\frac{1}{2}}
    \left(\sum\limits_{k = n + 1}^{\infty}  \frac{1}{k^{2m}}\right)^{\frac{1}{2}}
    \label{th472:uneq1}
  \end{gather}
  (В силу неравенства Коши-Буняковского). \\
  Так как
  \begin{gather*}
    \left(\sum\limits_{k = n + 1}^{\infty}
    \frac{1}{k^{2m}}\right)^{\frac{1}{2}} \leq \left(\int\limits_{n+1}^{\infty}
    \frac{dx}{x^{2m}} \right)^\frac{1}{2} = \frac{1}{\sqrt{2m-1}} \cdot
    \frac{1}{n^{m-1/2}}
  \end{gather*}
  То из \eqref{th472:uneq1} получаем \eqref{th472:eq1}
\end{proof}

\begin{remark}
  Поскольку $a_k(f) = C_k(f) + C_{-k}(f), \ b_k(f) = i(C_k(f) - C_{-k}(f))$, то
  из \eqref{th471:eq2} следует, что
  \begin{gather*}
    |a_k(f)| = \frac{\alpha_k}{k^m}, \ |b_k(f)| = \frac{\beta_k}{k^m}, \ k \in
    \mathbb{N}, \ \sum \alpha_k^2 < \infty, \ \sum \beta_k^2 < \infty
  \end{gather*}
\end{remark}


  \part{Интегралы}
  \chapter{Кратные интегралы}
  \section{Определение интеграла Римана на $n$-ом промежутке}
Пусть $\mathbb{R}^n$ --- $n$-мерное арифметическое евклидово пространство(ЕП).
\\ ($\mathbb{R}^2$ отождествляется $oxy$, $\mathbb{R}^3$ с $oxyz$).\\
$(x_1, \dots, x_n) \in \mathbb{R}^n$ далее обозначается через $x$.

\begin{definition}
  Множество $I = \{x \in \mathbb{R}^n : a_i \leq x \leq b_i, \ i = 1, \dots, n\}$
  называется промежутком или координатным параллелепипедом. По аналогии с
  одномерным случаем записывают:
  \begin{gather*}
    a = (a_1, \dots, a_n), \ b = (b_1, \dots, b_n)
  \end{gather*}
\end{definition}

\begin{definition}
  Число $|I| := \prod\limits_{i=1}^{n} (b_i - a_i)$ называют объемом, либо мерой
  промежутка. \\
  Объем(меру) обозначают символами: $v(I), \mu(I)$.
  \begin{align*}
    \text{При} \ n &= 1, \mu(I) \text{ --- длина отрезка} \\
    \text{При} \ n &= 2, \mu(I) \text{ --- площадь прямоугольника} \\
    \text{При} \ n &= 3, \mu(I) \text{ --- объем прямоугольного параллелепипеда}
  \end{align*}
  Разобьем каждый из координатных отрезков $[a_i, b_i], \ i = 1, \dots, n$ на
  конечное число более мелких отрезков. Эти разбиения индуцируют разбиение
  промежутка $I$ на более мелкие промежутки, получающиеся прямым произведением
  промежутков $a_i, b_i$.
\end{definition}

\begin{definition}[Декартово произведение]
  $X \times Y := \{(x, y) : x \in X, y \in Y \}$
\end{definition}

\begin{definition}
  Описанное представление промежутка $I$ в виде объеденения промежутков $i, j$
  из более мелких промежутков $i, j$ назовем разбиением и обозначим $T(k =
  k_T)$.
\end{definition}

\begin{definition}
  $\lambda(T) = \max\{d(I_1), \dots, d(I_k)\}$.
  $d(I_j)$ называется диаметром разбиения~$T$. Пусть $f(x) = f(x_1, \dots,
  x_n)$ --- функция, определенная на промежутке $I, \\ T = \{I_1, \dots, I_k\}$
  --- разбиение промежутка $I$. \\
  $\xi = (\xi^1, \dots, \xi^k)$ --- некоторый набор точек $\xi^j$, таких что
  $\xi^j \in I_j$.
\end{definition}

\begin{definition}
  Сумма $\sigma(f, T, \xi) := \sum\limits_{j = 1}^k f(\xi^j)|I_j|$ называется
  интегральной суммой Римана.
\end{definition}

\begin{definition}
  Если существует конечный предел $\mathcal{J}$, $\sigma(f, T, \xi), \
  \lambda(T) \to 0$, то его называют интегралом Римана от функции $f$ на
  промежутке $I$:
  \begin{gather}
    \mathcal{J} = \int\limits_I f(x) dx := \lim\limits_{\lambda(T) \to 0}
    \sigma(f, T, \xi)
    \label{def517:eq1}
  \end{gather}
  Функцию $f$ называют интегрируемой на промежутке $I$. Множество таких функций
  обозначим $\mathcal{R}(I)$. Уточним, что равенство \eqref{def517:eq1}
  означает:
  $\forall \varepsilon > 0 \ \exists \delta > 0 : \forall T = \{I_1, \dots,
  I_k\}, \lambda(T) < \delta$ и при любом выборе точек $\xi^j \in I_j, j = 1,
  \dots, k$ выполняется неравенство:
  \begin{gather*}
    \left|\mathcal{J} - \sum\limits_{j = 1}^{k} f(\xi^j) |I_j|\right| <
    \varepsilon
  \end{gather*}
  Равносильные отображения интеграла таковы:
  \begin{gather*}
    \int\limits_I f(x_1, \dots, x_n) \ dx_1 \dots dx_n, \
    \underset{n}{\underbrace{\int \dots \int}} f(x_1, \dots, x_n) \ dx_1 \dots
    dx_n, \ n \in \mathbb{N}
  \end{gather*}
\end{definition}

Мы видим, что данное определение повторяет определение интеграла Римана, и при
$m = 1$ совпадает. Схожесть опеределений позволяет найти сходные методы о
решении вопроса условий существования.

\section{Условие существования кратного интеграла}

\begin{theorem}
  Если $f \in \mathcal{R}(I)$, то $f$ ограниченна на $I$.
  Пусть функция $f$ определена на $I, \ I \in \mathbb{R}^n$. $T = \{I_j\}$ ---
  разбиение промежутка $I$.
  \begin{gather*}
    m_j = \inf\limits_{x \in I_j} f(x), \ M_j = \sup\limits_{x \in I_j} f(x)
  \end{gather*}
\end{theorem}

\begin{definition}
  Величины $s(f, T) = \sum\limits_{j} m_j |I_j|, \ S(f,T) = \sum\limits_{j} M_j
  |I_j|$ называются соотвтетственно нижней и верхней суммой Дарбу на промежутке
  отвечающему разбиению $T$ этого промежутка.
\end{definition}

\begin{comment}
  Совершенно аналогично доказывается при $m = 1$ теорема \eqref{th522}.
\end{comment}

\begin{theorem}
  \label{th522}
  Для того, чтобы ограниченная на промежутке функция $f$ была интегрируема на
  $I \Longleftrightarrow$
  \begin{gather}
    \lim\limits_{\lambda(T) \to 0} (S(f, T) - s(f, T)) = 0
    \label{th522:eq1}
  \end{gather}
\end{theorem}

\begin{remark}[критерий интегрируемости Римана]
  Если обозначить колебания $M_j - m_j, \ I_j$ через $\Omega(f, I_j)$, то
  \eqref{th522:eq1} можно записать в виде:
  \begin{gather*}
    \lim\limits_{\lambda(T) \to 0} \sum\limits_{j} \Omega(f, I_j) |I_j| = 0
  \end{gather*}
\end{remark}

\begin{definition}
  Говорят, что множество $\mathbb{E} \ n$-мерного пространства имеет меру $0$ в
  смысле Жордана или имеет нулевой $n$-мерный объем, если для $\forall
  \varepsilon > 0$ существует покрытие множества $\mathbb{E}$ конечной системой
  $\{I_j\} \ n$-мерных промежутков. $\sum\limits_{j} |I_j|$ объемов которых
  меньше~$\varepsilon$. В этом случае пишем: $\mu(\mathbb{E}) = 0$
\end{definition}

\begin{theorem}
  \label{th523}
  Пусть функция $f$ ограничена на $n$-мерном промежутке $I \subset
  \mathbb{R}^n$ и $\mathbb{E}_f$ --- множество ее точек разрыва. Тогда, если
  $\mu(\mathbb{E}_f) = 0$, то $f \in \mathcal{R}(I)$ (интегрируема).
\end{theorem}

\begin{comment}
  Эта теорема при $n = 1$ обычно устанавливается в разделе ``определенный
  интеграл''. В общем случае доказывается аналогично.
\end{comment}

\section{Кратный интеграл по множеству}
Пусть функция $f$ определена на $\mathbb{E} \subset \mathbb{R}^n$, условимся
символом $f_{\mathbb{E}}(x)$ обозначать функцию $= 0$ вне $\mathbb{E}$.
\begin{definition}
  \label{def531}
  Интеграл функции по множеству $\mathbb{E}$ определяется соотношением:
  \begin{gather}
    \int\limits_{\mathbb{E}} f(x) dx := \int\limits_I f_{\mathbb{E}}(x) dx
    \label{def531:eq1}
  \end{gather}
  где $I$ --- наименьший промежуток, содержащий $\mathbb{E}$. \\
  Если стоящий справа интеграл существует, то он называется интегрируемым по
  Риману на~$\mathbb{E}$.
  Совокупность всех интегрируемых на множестве $\mathbb{E}$ функций обозначим
  $\mathcal{R}(\mathbb{E})$.
\end{definition}

\begin{definition}
  Множество $\mathbb{E} \subset \mathbb{R}^n$ назовем допустимым, если оно
  ограничено в $\mathbb{R}^n$ и его граница $\partial \mathbb{E}$ есть
  множество меры нуль (в смысле Жордана).
\end{definition}

\begin{example}
  Куб, тетраэдер, шар и т. д. являются допустимыми множествами.
\end{example}

\begin{comment}
  Граница $\partial \mathbb{E}$ множества $\mathbb{E}$ состоит из точек, в
  любой окрестности которых имеются как точки из $\mathbb{E}$, так и точки из
  дополнения $\mathbb{E}$ до $\mathbb{R}^n$.
\end{comment}

\begin{theorem}
  Пусть $\mathbb{E}$ --- допустимое множество в $\mathbb{R}^n$. $f$ ---
  функция, определенная на $\mathbb{E}$ и пусть множество точек разрыва
  $\mathbb{E}_f$ множества $\mathbb{E}$ имеют нулевую меру, тогда функция $f
  \in \mathcal{R}(\mathbb{E})$ (интегрируема).
\end{theorem}

\begin{proof}
  Функция $f_{\mathbb{E}}$ по сравнению с функцией $f$ может иметь
  дополнительные разрывы на $\partial \mathbb{E}$, которые, по условию, имеет
  меру нуль. Поэтому множество точек разрыва функции $f_{\mathbb{E}}$ так же
  имеет нулевую меру. \\
  Отсюда из определения \eqref{def531} и теоремы \eqref{th523} следует, что
  функция интегрируема.
\end{proof}

\section{Мера(объем) множества}
\begin{definition}
  Мерой Жордана(или объемом) ограниченного множества $\mathbb{E}$ назовем
  величину $\mu(\mathbb{E}) := \int\limits_{\mathbb{E}} dx$, если указанный
  интеграл(Римана) существует. В последнем случае множество $\mathbb{E}$
  называют измеримым в смысле Жордана.
\end{definition}

\begin{theorem}
  Допустимое множество измеримо в смысле Жордана.
\end{theorem}

\begin{proof}
  Рассмотрим характеристическую функцию:
  \begin{gather*}
    \aleph_{\mathbb{E}}(x) =
    \begin{cases}
      1, \ x \in \mathbb{E} \\
      0, \ x \not \in \mathbb{E}
    \end{cases}
  \end{gather*}
  \begin{comment}
    $\aleph$ --- Алеф.
  \end{comment}
  Очевидно, функция $\aleph_{\mathbb{E}}(x)$ имеют разрывы в граничных и только
  в граничных точках множества $\mathbb{E}$. \\
  По определению \eqref{def531}:
  \begin{gather}
    \int\limits_{\mathbb{E}} 1 \cdot dx = \int\limits_I \aleph_{\mathbb{E}}(x)
    dx
    \label{th541:eq1}
  \end{gather}
  где $I$ --- наименьший промежуток, содержащий множество $\mathbb{E}$. \\
  А так как множество точек разрыва $\aleph_{\mathbb{E}}(x)$ совпадает с
  границей $\partial \mathbb{E}$ и $\mu(\partial \mathbb{E}) = 0$, то по
  теореме \eqref{th523} интеграл \eqref{th541:eq1} существует. \\
  $\mu(\mathbb{E})$ носит комплексный смысл, если $\mathbb{E}$ --- измеримое
  множество, и
  \begin{gather*}
    \int\limits_I \aleph_{\mathbb{E}}(x) dx = \lim\limits_{\lambda(T) \to 0}
    s(\aleph_{\mathbb{E}}, T) = \lim\limits_{\lambda(T) \to 0}
    S(\aleph_{\mathbb{E}}, T), \ (I \supset \mathbb{E})
  \end{gather*}
  то
  \begin{gather*}
    s(\aleph_{\mathbb{E}}, T) = \sum\limits_{j} m_j |I_j|, \
    S(\aleph_{\mathbb{E}}, T) = \sum\limits_{j} M_j |I_j|
  \end{gather*}
  где $s, S$ --- нижняя и верхняя суммы Дарбу. Но в силу определения функции
  $\aleph_{\mathbb{E}}: s(\aleph_{\mathbb{E}}, T)$ равна сумме объемов
  промежутков $I_j$, лежащих в множестве $\mathbb{E}$ (это объем вписанного в
  $\mathbb{E}$ многогранника), а $S(\aleph_{\mathbb{E}}, T)$ равна сумме
  объемов тех промежутков $I_j$, которые имеют общие точки с множеством
  $\mathbb{E}$ (Объем описанного многогранника).
\end{proof}

\begin{approval}
  Мера $\mu(\mathbb{E})$ есть общий предел, при $\lambda(T) \to 0$ объемов,
  вписанных в $\mathbb{E}$ и описанных около $\mathbb{E}$ многогранников.
\end{approval}

\begin{remark}
  Можно показать, что измеримы по Жордану только измеримые множества
  $\Longleftrightarrow$ множество $\mathbb{E}$ является измеримым по Жордану
  $\Longleftrightarrow$ его границы имеют меру нуль в смысле Жордана.
\end{remark}

\begin{remark}
  При $n = 2$ понятие измеримого по Жордану множества совпадает с понятием
  квадрируемой плоской фигуры и меры Жордана с ее площадью.
\end{remark}

\section{Свойства кратных интегралов}
На кратные интегралы по ограниченной функции переносятся все свойства
интегралов по отрезку. \\
Доказатаельства аналогичны одномерному случаю:
\begin{enumerate}
  \item \underline{Линейность} интеграла по множеству. \\
    $f_1, \dots, f_n$ интегрируема на множестве $\mathbb{E}, \ \lambda_1, \dots, \lambda_n$
    \\
    $\lambda_1 f_1, \dots, \lambda_n f_n$ так же интегрируема на $\mathbb{E}$ и
    \begin{gather*}
      \int\limits_{\mathbb{E}} \left(\lambda_1 f_1(x) + \dots + \lambda_n
      f_n(x) \right) dx = \lambda_1 \int\limits_{\mathbb{E}} f_1(x) dx + \dots
      \lambda_n \int\limits_\mathbb{E} f_n(x) dx
    \end{gather*}
  \item \underline{Аддитивность} интеграла по множеству. \\
    Если $\mathbb{E}_1, \mathbb{E}_2$ --- допустимые множества в $\mathbb{R}^n$
    и $\mu(\mathbb{E}_1 \cap \mathbb{E}_2) = 0 \ (\mathbb{E}_1 \cap \mathbb{E}_2
    = 0$ в частности), а $f$ --- функция, определенная на $\mathbb{E}_1
    \cup \mathbb{E}_2$, то при условии существования интегралов имеет место
    равенство:
    \begin{gather*}
      \int\limits_{\mathbb{E}_1 \cup \mathbb{E}_2} f(x) dx =
      \int\limits_{\mathbb{E}_1} f(x) dx + \int\limits_{\mathbb{E}_2} f(x) dx
    \end{gather*}
  \item \underline{Общая оценка}. \\
    Если $f \in \mathcal{R}(\mathbb{E})$, то $|f| \in \mathcal{R}(\mathbb{E})$ и имеет место
    неравенство:
    \begin{gather*}
      \left| \int\limits_{\mathbb{E}} f(x) dx \right| \leq
      \int\limits_{\mathbb{E}} |f(x)| dx
    \end{gather*}
  \item \underline{Интегрирование неравенств}.
    Если функции $f, g \in \mathcal{R}(\mathbb{E})$ и $f(x) \leq g(x), \ x \in
    \mathbb{E}$, то:
    \begin{gather*}
      \int\limits_\mathbb{E} f(x) dx \leq \int\limits_\mathbb{E} g(x) dx
    \end{gather*}
  \item \underline{Следствие из 4}. \\
    Если $f$ интегрируема и $m \leq f(x) \leq M, x \in \mathbb{E}$
    \begin{gather*}
      m\mu(\mathbb{E}) \leq \int\limits_\mathbb{E} f(x)dx \leq M\mu(\mathbb{E})
    \end{gather*}
  \item \underline{Теорема о среднем}. \\
    Если в дополнении условия 5. множество $\mathbb{E}$ линейно связано, а $f$
    --- непрерывна на $\mathbb{E}$, то существует $\xi \in \mathbb{E}$, такое
    что:
    \begin{gather*}
      \int f(x) dx = f(\xi) \mu(\mathbb{E})
    \end{gather*}
\end{enumerate}

\section{Сведение кратного инетграла к повторному}
\begin{theorem}
  $X \times Y \subset \mathbb{R}^{m+n}$ является прямым произведением
  промежутков: $X \subset \mathbb{R}^m$ и $Y \subset \mathbb{R}^n$. \\
  Если для $f(x, y)$ определенной на $X \times Y, (x \in X, y \in Y)$
  существует интеграл:
  \begin{gather}
    \int\limits_{X \times Y} f(x, y) dx dy
    \label{th561:ex1}
  \end{gather}
  и для любого $x \in X$ существует:
  \begin{gather}
    \mathcal{J}(x) = \int\limits_Y f(x, y) dy
    \label{th561:eq1}
  \end{gather}
  то существует так же и повторный интеграл:
  \begin{gather}
    \int\limits_X dx \int\limits_Y f(x, y) dy
    \label{th561:ex2}
  \end{gather}
  и выполняется равенство:
  \begin{gather}
    \int\limits_{X \times Y} f(x, y) dx dy = \int\limits_X dx \int\limits_Y
    f(x, y) dy
    \label{th561:eq2}
  \end{gather}
\end{theorem}

\begin{proof}
  Любое разбиение $T$ промежутка $X \times Y$ индуцируется собственными
  разбиениями $T_x, T_y$. При этом, каждый промежуток разбиения $T$ есть прямое
  произведения $X_i \times Y_j, \ X_i, Y_j$ --- разбиения $T_x, T_y$. \\
  Очевидно, $|X_i \times Y_j| = |X_i| \cdot |Y_j|$. \\
  Положим $m_{ij} = \inf\limits_{x \in X_i, \ y \in Y_j} f(x, y), \ M_{ij} =
  \sup\limits_{x \in X_i, \ y \in Y_j} f(x, y)$, так что выполняется следующее
  неравенство:
  \begin{gather}
    m_{ij} \leq f(x, y) \leq M_{ij}, \ \forall (x, y) \in X_i \times Y_j
    \label{th561:uneq1}
  \end{gather}
  Фиксируем $\forall x \in X_i: x = \xi_i$. \\
  Учитывая \eqref{th561:uneq1} получим:
  \begin{gather*}
    m_{ij} |Y_j| \leq \int\limits_{Y_j} f(\xi_i, y) dy \leq M_{ij} |Y_j|
  \end{gather*}
  Просуммировав все по $j$ получим:
  \begin{gather*}
    \sum\limits_{j} m_{ij} |Y_j| \leq \mathcal{J}(\xi_i) = \int\limits_Y
    f(\xi_i, y) dy \leq \sum\limits_j M_{ij} |Y_j|
  \end{gather*}
  Отсюда имеем:
  \begin{gather}
    \sum\limits_{i} |X_i| \sum\limits_j m_{ij} |Y_j| \leq \sum\limits_{i}
    \mathcal{J}(\xi_i) |X_i| \leq \sum\limits_{i} |X_i| \sum\limits_{j} M_{ij}
    |Y_j|
    \label{th561:uneq2}
  \end{gather}
  По середине мы получили интегральную сумму для функции $\mathcal{J}(x)$.
  Крайние челны --- это суммы Дарбу $s(f, T), S(f, T)$ для кратного интеграла
  \eqref{th561:ex1}. \\
  Например $\sum\limits_{i, j} m_{ij} |X_i \times Y_j|$. Таким образом
  неравенство \eqref{th561:uneq2} перепишется в виде:
  \begin{gather}
    s(f, T) \leq \sum\limits_{i} \mathcal{J}(\xi_i) |X_i| \leq S(f, T)
    \label{th561:uneq3}
  \end{gather}
  Так как кратный интеграл \eqref{th561:ex1} существует по условию, то при
  $\lambda(T) \to 0$ обе суммы Дарбу неравенства \eqref{th561:uneq3} стремятся
  к этому интегралу, откуда
  \begin{gather*}
    \left|\lim\limits_{\lambda(T) \to 0} \sum\limits_{i} \mathcal{J}(\xi_i)
    |X_i| \right| = \int\limits_{X \times Y} f(x, y) dx dy
  \end{gather*}
  Левая часть этого равенства есть повторный интеграл:
  \begin{gather*}
    \int\limits_X \mathcal{J}(x) dx = \int\limits_X dx \int\limits_Y f(x, y) dy
  \end{gather*}
\end{proof}

\begin{remark}
  Применяя эту теорему несколько раз, можно свести вычисление по $k$-мерному
  промежутку к вычислению $k$ одномерных интегралов.
\end{remark}

\begin{consequence}
  Пусть $\mathbb{D} \subset Oxy$ --- область, ограниченная двумя кривыми
  $\mathcal{J} = \varphi(x), \mathcal{J} = \psi(x), (\varphi(x) \leq \psi(x))$
  и двумя прямыми $x = a$ и $x = b$. \\
  Тогда, если для $f(x, y)$ существует $\iint\limits_\mathbb{D} f(x, y)
  dx dy, \ x \in [a,b]$. \\
  $\mathcal{J}(x) = \int\limits_{\varphi(x)}^{\psi(x)} f(x, y) dy$, то
  существует так же повторный интеграл: $\int\limits_a^b dx
  \int\limits_{\varphi(x)}^{\psi(x)} f(x, y) dy$ и выполняется равенство
  \begin{gather*}
    \iint\limits_\mathbb{D} f(x, y) dx dy = \int\limits_a^b dx
    \int\limits_{\varphi(x)}^{\psi(x)} f(x, y) dy.
  \end{gather*}
\end{consequence}

\section{Замена переменных в кратных интегралах}
Напомним, что обтображение:
\begin{gather*}
  x = \varphi(t) =
  \begin{cases}
    x_1 = \varphi_1(t_1, \dots, t_n) \\
    \dots \\
    x_n = \varphi_n(t_1, \dots, t_n)
  \end{cases}
\end{gather*}
называется регулярным в области $\mathbb{D} \subset \mathbb{R}^n$, если:
\begin{enumerate}
  \item $\varphi_1, \dots, \varphi_n$ имеют в $\mathbb{D}$ непрерывные частные
    производные по всем аргументам.
  \item Матрица Якоби:

    \begin{gather}
      \varphi'(t) :=
      \begin{pmatrix}
        \frac{\partial \varphi_1}{\partial t_1} & \dots & \frac{\partial
        \varphi_1}{\partial t_n} \\
        \vdots & \ddots & \vdots \\
        \frac{\partial \varphi_n}{\partial t_1} & \dots & \frac{\partial
        \varphi_n}{\partial t_n}
      \end{pmatrix}
      \label{ch57:mat1}
    \end{gather}
    Определитель(якобиан):
    \begin{gather*}
      \det \varphi'(t) =
      \left|
      \begin{matrix}
        \frac{\partial \varphi_1}{\partial t_1} & \dots & \frac{\partial
        \varphi_1}{\partial t_n} \\
        \vdots & \ddots & \vdots \\
        \frac{\partial \varphi_n}{\partial t_1} & \dots & \frac{\partial
        \varphi_n}{\partial t_n}
      \end{matrix}
      \right|
      =: \frac{\mathbb{D}(\varphi_1, \dots, \varphi_n)}{\mathbb{D}(t_1, \dots,
      t_n)}
    \end{gather*}
    Отличен от нуля в $\mathbb{D}$.
    Матрица Якоби \eqref{ch57:mat1} называется производной отображения
    $\varphi$, а линейный оператор в $\mathbb{R}^n$ называется диффиренциалом
    $\varphi$ в точке $t$, в области $\mathbb{D}, \ (\mathbb{D}_\varphi(t))$.
    Регулярное отображение является локально обратимым.
\end{enumerate}
\begin{lemma}
  Пусть $\varphi$ --- регулярное отображение области $\mathbb{D} \subset
  \mathbb{R}_t^n$, $I$ --- замкнутый промежуток, лежащий в $\mathbb{D}, \ (I
  \subset \mathbb{D})$. $\varphi(I)$ --- образ промежутка $I$, $(\varphi(I)
  \subset R_x^n)$. Тогда, существует такая точка $\tau \in I$, что:
  \begin{gather*}
    \mu(\varphi(I)) = |\det \varphi'(\tau)| \cdot |I|
  \end{gather*}
  \begin{comment}
    $|I| = \mu(I)$
  \end{comment}
\end{lemma}

\begin{clarification}
  Лемма при $n = 1$ следует из формулы Лагранжа:
  $\varphi(b) - \varphi(a) = \varphi(\tau)(b-a)$, \\ где $\tau \in [a,b] = I$,
  поскольку при $\varphi'(\tau) \not = 0$ функция $\varphi$ монотонна и
  \begin{gather*}
    |\varphi(b) - \varphi(a)| = \mu(\varphi(I))
  \end{gather*}
  Если $n \geq 2$ и отображение $\varphi$ является линейным преобразованием
  \begin{gather*}
    x = \varphi(t) =
    \begin{cases}
      x_1 = a_{11}t_1 + \dots + a_{1n} t_n \\
      \dots \\
      x_n = a_{n1}t_1 + \dots + a_{nn} t_n \\
    \end{cases}
  \end{gather*}
  с матрицей $A = (a_{ij}) = \varphi'(t), \ x = At$, то образ промежутка
  $\varphi(I)$ является параллелепипедом $\varphi(I) \subset \mathbb{R}_x^n$,
  объем которого равен:
  \begin{gather*}
    |\det A| \cdot |I| = |\det \varphi'| \cdot |I|
  \end{gather*}
  При $n = 3, I$ --- прямоугольный параллелепипед, построенный на векторах:
  $h^1 = (h_1, 0, 0), h^2 = (0, h_2, 0), h^3 = (0, 0, h_3)$, тогда
  $\varphi(h^i) = (a_{1i} h_i, a_{2i} h_i, a_{3i} h_i), \ i = 1, 2, 3$. Отсюда
  объем параллелепипеда, построенного на векторах $\varphi(h^1), \varphi(h^2),
  \varphi(h^3)$ как на ребрах, равен
  \begin{gather*}
    |<\varphi(h^1), \varphi(h^2), \varphi(h^3)>| = |\det A| \cdot |h_1 h_2 h_3|
  \end{gather*}
  Рассматривая уже общий, нелинейный случай, следует принять во внимание, что в
  малой окрестности точки $t \in \mathbb{D}$ является почти линейным
  отображением:
  \begin{gather*}
    \varphi(t + h) = \varphi(t) + \mathbb{D}_\varphi(t) +
    \bar{\bar{o}}\left(h\right), \ \text{при} \ h \to 0
  \end{gather*}
  Поэтому, если размеры промежутка $I$ малы, то с малой относительной
  погрешностью можно сказать:
  \begin{gather*}
    \mu(\varphi(I)) \approx |\det \varphi'(t)| \cdot |I|, \ t \in I
  \end{gather*}
  В такой ситуации и используется лемма ($I$ --- малый промежуток).
\end{clarification}

\begin{theorem}
  Пусть $\varphi: \mathbb{E}_t \to \mathbb{E}_x$ --- отображение измеримого (по
  Жордану) множества $\mathbb{E}_t \subset \mathbb{R}_t^n, \ \mathbb{E}_x
  \subset \mathbb{R}_x^n$, при чем отображение $\varphi$ регулярно в некоторой
  области $\mathbb{D}$, содержащей замыкание $\overline{\mathbb{E}}_t,
  \mathbb{E}_t, (\overline{\mathbb{E}}_t \subset \mathbb{D})$. Тогда, если
  $f(x) \in \mathcal{R}(\mathbb{E}_x)$, то $f(\varphi(t)) |\det \varphi'(t)|
  \in \mathcal{R}(\mathbb{E}_t)$ и имеет место равенство:
  \begin{gather}
    \int\limits_{\mathbb{E}_x} f(x) dx = \int\limits_{\mathbb{E}_t}
    f(\varphi(t)) |\det \varphi'(t)| dt
    \label{th571:eq1}
  \end{gather}
\end{theorem}

\begin{proof}
  Случай, когда $\mathbb{E}_t$ --- промежуток, а $f(x)$ ограничена и непрерывна
  на $\mathbb{E}_x$. В этом случае функция $g(t) = f(\varphi(t))|\det
  \varphi'(t)|$ так же ограничена и непрерывна на $\mathbb{E}_t$, а
  следовательно и интегрируема на $\mathbb{E}_t$. Любому разбиению $T$
  промежутка $\mathbb{E}_t$ на промежутке $I_1, \dots, I_k$ соответствует
  разложение множества $\varphi(\mathbb{E}_t) = \mathbb{E}_x$ на множество
  $\varphi(I_j), \ j = 1, \dots, k$. \\
  Все эти множества измеримы, связаны и пересекаются попарно лишь по множествам
  меры нуль. Поэтому, в силу аддитивности интеграла:
  \begin{gather}
    \int\limits_{\mathbb{E}_x} f(x) dx = \sum\limits_{j = 1}^{k}
    \int\limits_{\varphi(I)} f(x) dx
    \label{th571:eq2}
  \end{gather}
  По теореме о среднем:
  \begin{gather}
    \int\limits_{\varphi(I_j)} f(x) dx = f(\xi^j) \mu(\varphi(I_j)), \ \xi^j
    \in \varphi(I_j)
    \label{th571:eq3}
  \end{gather}
  Пусть $\eta^j = \varphi^{-1}(\xi^j) \in I_j$, так что $\xi^j =
  \varphi(\eta^j)$. \\
  Поскольку по лемме $\mu(\varphi(I_j)) = |\det \varphi'(\tau^j)|\cdot |I_j|$,
  где $\tau^j \in I_j$, то из \eqref{th571:eq2} и \eqref{th571:eq3} получаем:
  \begin{gather}
    \int\limits_{\mathbb{E}_x} f(x) dx = \sum\limits_{j = 1}^{k}
    f(\varphi(\eta^j))\left|\det \varphi'(\tau^j)\right|\cdot |I_j| =: \sigma_1
    \label{th571:eq4}
  \end{gather}
  Составим сумму: $\sigma_2 := \sum\limits_{j = 1}^{k} f(\varphi(\eta^j))
  \left|\det \varphi'(\eta^j)\right| \cdot |I_j|$, которая является
  интегрируемой суммой для функции $g(t), \ t \in \mathbb{E}_t$. Поскольку
  $g(t) \in \mathcal{R}$, то:
  \begin{gather}
    \lim\limits_{\lambda(T) \to 0} \sigma_2 = \int\limits_{\mathbb{E}_t} g(t)
    dt
    \label{th571:eq5}
  \end{gather}
  Положим $\psi(t) = |\det \varphi'(t)|$ и оценим разность $\sigma_1 -
  \sigma_2$:
  \begin{gather*}
    \sigma_1 - \sigma_2 = \sum\limits_{j = 1}^{k} f(\xi^j) (\psi(\tau^j) -
    \psi(\eta^j)) |I_j|
  \end{gather*}
  Функция $\psi(t)$ непрерывна на замкнутом множестве $\mathbb{E}_t$ и по
  теореме Кантора равномерно непрерывна на $\mathbb{E}_t$, так что $\forall
  \varepsilon > 0 \ \exists \delta > 0, \ \lambda(T) < \delta$:
  \begin{gather*}
    |\psi(\tau^j) - \psi(\eta^j)| < \varepsilon, \ \forall j \\
    \rho(\tau^j, \eta^j) < \delta
  \end{gather*}
  Отсюда, при $\lambda(T) < \delta$:
  \begin{gather*}
    |\sigma_1 - \sigma_2| \leq \max\limits_{x \in \mathbb{E}_x} |f(x)| \cdot
    \varepsilon \sum\limits_{j = 1}^{k} |I_j| = C \varepsilon
  \end{gather*}
  где $C$ не зависит от $\varepsilon$, $C = const$.
  \begin{gather}
    \lim\limits_{\lambda(T) \to 0} (\sigma_1 - \sigma_2) = 0
    \label{th571:eq6}
  \end{gather}
  Равенство \eqref{th571:eq1} теперь следует из \eqref{th571:eq4},
  \eqref{th571:eq5} и \eqref{th571:eq6}. \\
  Доказательство теоремы в общем случае можно провести, придерживаясь
  рассмотренной схемы.
\end{proof}

\begin{consequence}
  Величина интеграла от $f$ по множеству $\mathbb{E} \subset \mathbb{R}^n$ и
  зависит от выбора декартовых координат в $\mathbb{R}^n$.
\end{consequence}

\begin{proof}
  Пусть $\mathbb{E}_x, \mathbb{E}_t$ --- запись множества $\mathbb{E}$.
  $p$ --- точка множества $\mathbb{E}, \ x = (x_1, \dots, x_n)$ --- ее
  координаты в первой системе, $t = (t_1, \dots, t_n)$ --- во второй системе.
  Тогда $f(p) = f_x(x_1, \dots, x_n) = f_t(t_1, \dots, t_n)$, где $f_t = f_x
  \circ \varphi$ (суперпозиция). Поскольку переход от одной системы декартовых
  координат к другой имеет якобиан по модулю равный единице, то есть:
  \begin{gather*}
    \int\limits_{\mathbb{E}_x} f_x(x) dx = \int\limits_{\mathbb{E}_t}
    f_x(\varphi(t))\left|\det \varphi'(t)\right| dt = \int\limits_{\mathbb{E}_t} f_t(t) dt
  \end{gather*}
\end{proof}


  \chapter{Криволинейные интегралы}
  \section{Криволинейный интеграл первого рода}
Пусть $\Gamma = \{M(S) : 0 \leq S \leq S_0\}$ определенная кривая в
$\mathbb{R}^3$, в частности $\mathbb{R}^2. \\ M(S) = (x(S), y(S), z(S))$ --- ее
параметрическое представление, где в качестве переменной взята длина дуги $S$.
\\
Кривая, у которой переменная длина дуги отсчитывается от точки $A$
обозначим: $\Gamma = \wideparen{AB}$. Противоположно ориентированную прямую, у
которой переменная длина дуги отсчитывается от точки $B$ обозначим $\Gamma =
\wideparen{BA}$.

\begin{definition}
  Пусть на токах $M(S)$, кривой $\Gamma = \wideparen{AB}$ задана некотороая
  функция $F$. Выражение $\int\limits_\Gamma F(x, y, z) dS$ опеределенное по
  формуле:
  \begin{gather}
    \int\limits_\Gamma F(x, y, z) dS = \int\limits_0^{S_0} F(x(S), y(S), z(S))
    dS
    \label{def611:eq1}
  \end{gather}
  называется криволинейным интегралом первого рода (КИПР) от функции $F$ по
  кривой $\Gamma$. Этот интеграл обозначается символами:
  \begin{gather*}
    \int\limits_\Gamma F(M(S)) dS, \ \int\limits_\Gamma F dS, \
    \int\limits_{\wideparen{AB}} FdS
  \end{gather*}
\end{definition}

Отметим свойства интеграла \eqref{def611:eq1}:
\begin{enumerate}
  \item $\int\limits_\Gamma dS = S_0$ --- длина кривой $\Gamma$.
  \item Если $F$ непрерывна на $\Gamma$ (то есть $F(x(S), y(S), z(S))$
    непрерывная на отрезке $S \in [0, S_0]$), то $\int\limits_\Gamma FdS$
    существует.
  \item КИПР не зависит от ориентации кривой:
    \begin{gather*}
      \int\limits_{\wideparen{BA}} FdS = \int\limits_{\wideparen{AB}} FdS
    \end{gather*}
\end{enumerate}

\begin{definition}
  Кривую $\Gamma$ называют гладкой, если в ее параметрическом представлении:
  \begin{gather}
    x = \varphi(t), \ y = \psi(t), \ z = \chi(t), \ a \leq t \leq b
    \label{def612:eq2}
  \end{gather}
  все функции непрерывно диффиренцируемы на отрезке $[a, b]$ и $\varphi'(t)^2 +
  \psi'(t)^2 + \chi'(t)^2~>~0, \ \forall t \in~[a, b]$.
\end{definition}

\begin{theorem}
  Пусть $\Gamma$ --- гладкая кривая с параметрическим представлением
  \eqref{def612:eq2},\\ $F$ --- непрерывная на $\Gamma$ функция. Тогда:
  \begin{gather}
    \int\limits_\Gamma F(x, y, z) dS = \int\limits_a^b F(\varphi(t), \psi(t),
    \chi(t)) \sqrt{\varphi'(t)^2 + \psi'(t)^2 + \chi'(t)^2} dt
    \label{th611:eq1}
  \end{gather}
\end{theorem}

\begin{proof}
  Кривая $\Gamma$ спрямляема. Переменную дуги $S = S(t)$ можно принять в
  качестве параметра, и тогда $\int\limits_\Gamma FdS$ существует. Учитывая,
  что $S'(t) = \sqrt{\varphi'(t)^2 + \psi'(t)^2 + \chi'(t)^2}$ получим
  \eqref{th611:eq1}.
\end{proof}

\begin{remark}
  Если плоская кривая $\Gamma$ является графиком функци $y = f(x), a \leq x
  \leq b$, то ее представлением является функция $x = x, y = f(x)$. И тогда
  формула \eqref{th611:eq1} примет вид:
  \begin{gather*}
    \int\limits_\Gamma F(x, y) dS = \int\limits_a^b F(x, f(x)) \sqrt{1 +
    f'(x)^)} dx
  \end{gather*}
\end{remark}

\section{Криволинейный интеграл второго рода}

Пусть $\Gamma$ --- гладкая ориентированная кривая,\\ $\vec{r} = \vec{r}(S) =
\{x(S), y(S), z(S)\}, \ 0 \leq S \leq S_0$ --- ее векторное представление, в
котором за параметр $S$ взята переменная длина ее дуги. $\vec{\tau} =
\frac{d \vec{r}}{dS} = \{\frac{dx}{dS}, \frac{dy}{dS}, \frac{dz}{dS}\}$ ---
единичный касательный вектор к кривой $\Gamma$, его направление соответствует
выбранному направлению отсчета длины дуг. Если $\alpha, \beta, \gamma$ ---
углы, которые $\vec\tau$ образует с координатными осями, то $\vec\tau = \{\cos
\alpha, \cos \beta, \cos \gamma\}$. Получаем $\cos \alpha~=~\frac{dx}{dS}, \cos
\beta~=~\frac{dy}{dS}, \cos \gamma~=~\frac{dz}{dS}$. \\

Пусть на $\Gamma$ задана вектор-функция $\vec a = \vec a(x, y, z) = \vec
a(x(S), y(S), z(S)), \ 0 \leq S \leq S_0$. \\ Пусть $P, Q, R$ --- координаты
вектора $\vec a : \vec a = \{P, Q, R\}$. Функции $P, Q, R$ являются функциями
точки кривой $\Gamma$.

\begin{definition}
  \label{def621}
  $\int\limits_\Gamma \vec a d \vec r$ определенное по формуле:
  \begin{gather}
    \int\limits_\Gamma (\vec a, \vec r) dS
    \label{def621:int1}
  \end{gather}
  называется криволинейным интегралом второго рода (КИВР) по кривой $\Gamma =
  \wideparen{AB}$.
\end{definition}

Для интеграла $\int\limits_\Gamma \vec a d \vec r$ используется так же
обозначение:
\begin{gather*}
  \int\limits_\Gamma Pdx + Qdy + Rdz
\end{gather*}
где $Pdx + Qdy + Rdz$ --- записанное в координатной форме скалярное
произведение векторов $\vec a = \{P, Q, R\}$ и $d \vec r = \{dx, dy, dz\}$.
Учитывая, что $(\vec a, \vec \tau) = P \cos \alpha + Q \cos \beta + R \cos
\gamma$, то определение \eqref{def621} можно записать в виде:
\begin{gather}
  \int\limits_\Gamma \vec a d \vec r \equiv \int\limits_\Gamma Pdx + Qdy + Rdz
  := \int\limits_\Gamma (P\cos \alpha + Q \cos \beta + R \cos \gamma) dS
  \label{def621:int2}
\end{gather}

Следующие свойства КИВР под силу доказать каждому:
\begin{enumerate}
  \item $\vec a = \vec a(x, y, z), \ \Gamma = \wideparen{AB}$, то
    $\int\limits_\Gamma \vec a d \vec r$ существует.
  \item КИВР меняет знак при изменении ориентации кривой $\Gamma$:
    \begin{gather*}
      \int\limits_{\wideparen{BA}} \vec a d \vec r =
      -\int\limits_{\wideparen{AB}} \vec a d \vec r
    \end{gather*}
\end{enumerate}

\begin{theorem}
  Если $\Gamma = \wideparen{AB}$ --- гладкая ориентированная кривая. \\ $\vec r =
  \vec r(t) = \{\varphi(t), \psi(t), \chi(t)\}, \ a \leq t \leq b$ --- ее
  векторное представление ($\vec r(a) = A, \ \vec r(b) = B$), то получим:
  \begin{gather}
    \int\limits_\Gamma \vec a d \vec r = \int\limits_a^b (\vec a, \vec r') dt
    \label{th621:eq1}
  \end{gather}
\end{theorem}

\begin{proof}
  Поскольку
  \begin{gather}
    \vec \tau = \frac{d \vec r}{dS} = \frac{d \vec r}{dt} \cdot
    \frac{dt}{dS} = \frac{\vec r'}{S'}
    \label{th621:eq2}
  \end{gather}
  (здесь штрихом обозначены производные по $t$), то
  \begin{gather*}
    \int\limits_\Gamma \vec a d \vec r \overset{\eqref{def621:int1}} =
    \int\limits_\Gamma (\vec a, \vec \tau) dS \overset{\eqref{th611:eq1}} =
    \int\limits_a^b (\vec a, \vec \tau) S' dt \overset{\eqref{th621:eq2}} =
    \int\limits_a^b (\vec a, \vec r') dt
  \end{gather*}
\end{proof}

\begin{remark}
  В координатной форме формула \eqref{th621:eq1} примет вид:
  \begin{gather}
    \int\limits_\Gamma Pdx + Qdy + Rdz = \int\limits_a^b (P x' + Q y' + R z')dt
    \label{rem621:eq1}
  \end{gather}
\end{remark}

\begin{consequence}
  Если плоская кривая $\Gamma = \wideparen{AB}$ является графиком функции $y =
  f(x), \ a~\leq~x~\leq~b, \\ A = (a, f(a)), \ B = (b, f(b))$, то формула
  \eqref{th621:eq1} (при $x = t, \ R \equiv 0$) принимает вид:
  \begin{gather}
  \int\limits_\Gamma Pdx + Qdy + Rdz = \int\limits_a^b (P(x, f(x)) + Q(x,
  f(x))) f'(x) dx
  \label{cons621:eq1}
  \end{gather}
\end{consequence}

\begin{remark}
  Если $\vec a = \{P, 0, 0\}, \ Q \equiv R \equiv 0$, то $\int\limits_\Gamma
  \vec a d \vec r$ обозначается $\int\limits_\Gamma dx$, таким образом
  интеграл:
  \begin{gather}
    \int\limits_\Gamma P(x, y, z) dx = \int\limits_\Gamma P \cos \alpha \ dS
    \label{rem622:eq1}
  \end{gather}
  аналогично
  \begin{gather}
    \int\limits_\Gamma Q dy := \int\limits_\Gamma Q \cos \beta \ dS,
    \ \int\limits_\Gamma R dz := \int\limits_\Gamma R \cos \gamma \ dS
    \label{rem622:eq2}
  \end{gather}
  Отсюда, воспользовавшись адитивностью обычного интеграла, получим:
  \begin{gather*}
    \int\limits_\Gamma Pdx + Qdy + Rdz = \int\limits_\Gamma Pdx +
    \int\limits_\Gamma Qdy + \int\limits_\Gamma Rdz
    \label{rem622:eq3}
  \end{gather*}
\end{remark}

\begin{consequence}
  В условиях предыдущего следствия:
  \begin{gather}
    \int\limits_\Gamma Pdx = \int\limits_a^b P(x, f(x)) dx
    \label{cons621:eq2}
  \end{gather}
\end{consequence}

\begin{definition}
  Если кривая $\Gamma$ --- кусочно-гладкая, то есть представима в виде
  объединения конечного числа гладких ориентированных кривых $\Gamma_1, \dots,
  \Gamma_k, \ F, \ \vec a = \{P, Q, R\}$, определенный на точках кривой
  $\Gamma$, то пологают:
  \begin{gather*}
    \int\limits_\Gamma F dS := \sum\limits_{i = 1}^{k} \int\limits_{\Gamma_i} F
    dS, \int\limits_\Gamma \vec a d \vec r := \sum\limits_{i = 1}^{k}
    \int\limits_{\Gamma_i} \vec a d \vec r
  \end{gather*}
\end{definition}

\section{Формула Грина}
Пусть на плоскости $\mathbb{R}^2$ задана система координат по $Oxy$.

\begin{definition}
  Ориентация простого замкнутого контура $\Gamma$, лежащего на этой плоскости
  называется положительной, если она соответствует движению {\bfseries против}
  часовой стрелки. Противоположная ориентация называется отрицательной.
\end{definition}

Напомним, что простым замкнутым контуром в $\mathbb{R}^n$ называют $x =
\varphi(t) = (\varphi_1(t), \dots, \varphi_n(t)), \\ a \leq t \leq b$, у
которой нет других кратных точек.

\begin{definition}
  Пусть граница $\partial \mathbb{G}$, ограниченной плоской области
  $\mathbb{G}$,
  состоит из конечного числа простых замкнутых контуров. Совокупность этих
  контуров, ориентированных так, что при обходе каждого из них, область
  $\mathbb{G}$ остается слева(справа), называется положительной(отрицательной)
  ориентацией: $\partial \mathbb{G}^+ (\partial \mathbb{G}^-)$.
\end{definition}

\begin{definition}
  Грацница $\partial \mathbb{G}$ области $\mathbb{G}$ называется
  кусочно-гладкой, если она состоит из конечного числа простых кусочно-гладких
  контуров.
\end{definition}

\begin{remark}
  Если граница $\partial \mathbb{G}$ области $\mathbb{G}$ является кусочно
  гладкой, то ее площадь равна нулю, а само множество $\mathbb{G}$ квадрируемо.
\end{remark}

\begin{theorem}
  \label{th631}
  Если граница плоской, ограниченной области $\mathbb{G}$ является
  кусочно-гладкой, а функции $P, Q, \frac{\partial P}{\partial y},
  \frac{\partial Q}{\partial x}$ --- непрерывны на замыкании
  $\overline{\mathbb{G}}$ области $\mathbb{G}$, то:
  \begin{gather}
    \iint\limits_{\mathbb{G}} \left(\frac{\partial Q}{\partial x} -
    \frac{\partial P}{\partial y} \right) dx dy = \int\limits_{\partial
      \mathbb{G}^+} Pdx + Qdy
    \label{th631:eq1}
  \end{gather}
  эта формула называется {\bfseries формулой Грина}.
\end{theorem}

\begin{proof}
  Проведем его для произвольных областей.
\end{proof}

\begin{definition}
  Назовем область $\mathbb{G}$ областью $\mathbb{G}_y$, если $\mathbb{G}$ имеет
  вид:
  \begin{gather}
    \mathbb{G} = \{(x, y) : a < x < b, \ \varphi(x) < y < \psi(x)\}
    \label{def634:eq1}
  \end{gather}
  где $\varphi(x), \psi(x)$ --- кусочно-гладкие функции на $[a,b]$. Поменяв
  здесь $x$ и $y$ ролями получим определение области $\mathbb{G}_x$. Области,
  которые можно разрезать на куски вида $\mathbb{G}_x$ либо $\mathbb{G}_y$
  назовем простыми областями.
\end{definition}

\begin{lemma}
  Пусть в теореме \eqref{th631} $\mathbb{G}$ --- область по $\mathbb{G}_y$,
  тогда:
  \begin{gather}
    \iint\limits_{\mathbb{G}} \frac{\partial P}{\partial y} dx dy = -
    \int\limits_{\partial \mathbb{G}^+} Pdx
    \label{lem631:eq1}
  \end{gather}
\end{lemma}

\begin{proof}
  Сводя двойной интеграл к повторному, применяя формулу Ньютона-Лейбница и
  формулу \eqref{cons621:eq2} получим:
  \begin{gather}
    \iint\limits_{\mathbb{G}} \frac{\partial P}{\partial y} dx dy =
    \int\limits_a^b dx \int\limits_{\varphi(x)}^{\psi(x)} \frac{\partial
    P}{\partial y} dy = \int\limits_a^b \left[P(x, \psi(x)) - P(x,
    \varphi(x))\right] dx = \\
    -\int\limits_a^b P(x, \varphi(x)) dx - \int\limits_b^a P(x, \psi(x)) dx =
    -\int\limits_{\wideparen{AB}} P(x, y) dx - \int\limits_{\wideparen{CD}}
    P(x, y) dx
    \label{lem631:eq2}
  \end{gather}
  заметим, что $\int\limits_{\wideparen{BC}} Pdx = \int\limits_{\wideparen{DA}}
  Pdx = 0$, и что сумма: $\int\limits_{\wideparen{AB}} Pdx +
  \int\limits_{\wideparen{BC}} Pdx + \int\limits_{\wideparen{CD}} Pdx +
  \int\limits_{\wideparen{DA}} Pdx = \int\limits_{\partial \mathbb{G}^+} Pdx$.
  Поэтому из \eqref{lem631:eq2} получаем \eqref{lem631:eq1}. \\
  Аналогично, если $\mathbb{G}$ --- область типа $\mathbb{G}_x$, то:
  \begin{gather}
    \iint\limits_{\mathbb{G}} \frac{\partial Q}{\partial x} dx dy =
    \int\limits_{\partial \mathbb{G}^+} Qdy
    \label{lem631:eq3}
  \end{gather}
\end{proof}

\begin{lemma}
  Если в теореме \eqref{th631} область $\mathbb{G}$ допускает разбиение на
  области $\mathbb{G}_x, \mathbb{G}_y$, то выполняется \eqref{lem631:eq1}.
\end{lemma}

\begin{proof}
  Двойной интеграл в области $\mathbb{G}$ в силу адитивности есть сумма по
  кускам $\mathbb{G}_x, \mathbb{G}_y$, на которые разрезана $\mathbb{G}$. Для
  каждого справедливо \eqref{lem631:eq1}. Но соседние куски на общей части
  границ индуцируют противоположные ориентации. Поэтому, при сложении
  интегралов по границам всех кусков в результате взаимных уничтожений,
  останется $\int\limits_{\mathbb{G}} \partial \mathbb{G}$. \\
  Аналогично, если разрезается на $\mathbb{G}_x$, то справедливо
  \eqref{lem631:eq3}. Запищем для произвольной области выражения
  \eqref{lem631:eq1} и \eqref{lem631:eq3}, умножив \eqref{lem631:eq1} на $-1$.
  После сложения полученных соотношений, пользуясь свойством адитивности
  интеграла, получим формулу Грина.
\end{proof}

\begin{consequence}
  \begin{gather*}
    \mu(\mathbb{G}) = \iint\limits_{\mathbb{G}} dx dy \overset{(\rom{1})} =
    \int\limits_{\partial \mathbb{G}^+} x dy \overset{(\rom{2})} =
    -\int\limits_{\partial \mathbb{G}^+} y dx \overset{(\rom{3})} =
    \frac{1}{2} \int\limits_{\partial \mathbb{G}^+} xdy - ydx
  \end{gather*}
  \begin{align*}
    (\rom{1}) &\ \text{---} \ P = 0, Q = x \\
    (\rom{2}) &\ \text{---} \ P = -y, Q = 0 \\
    (\rom{3}) &\ \text{--- если возьмем полусумму}
  \end{align*}
\end{consequence}



\end{document}

